% !TeX root = ./../main.tex

\chapter{Zaplecze teoretyczne}\label{chapter:theoretical_background}

\agentinfo{
Rozdział wprowadza pojęcia, formalizm i definicje, z których korzystają rozdziały 4--9.
Docelowa objętość rozdziału: około 24 stron, więc każdy podrozdział musi być zwięzły.
Zasada podziału między zapleczem teoretycznym a przeglądem literatury:
tutaj znajdują się definicje, wzory i ustalone podstawy (wiedza podręcznikowa i definicje metod bazowych);
w przeglądzie literatury \agentref{chapters/literature_review.tex} znajduje się stan badań: kto, co, kiedy i z jakim wynikiem, porównania oraz luka badawcza.
Każde pojęcie jest definiowane raz, w tym rozdziale, a dalsze rozdziały odwołują się do niego przez etykiety.

Kolejność podrozdziałów odpowiada osi pracy: obliczenia kwantowe, szum, metody przeciwdziałania błędom, próbkowanie kwantowe, reprezentacje obrazów, uczenie maszynowe.

Źródła własne, z których można czerpać treść:
praca magisterska autora w języku polskim (podstawy informatyki kwantowej, błędy kwantowe, QCoSamp, PDU) \agentref{artykuły/magisterka_werner_final/chapters/01.tex} \agentref{artykuły/magisterka_werner_final/chapters/02.tex},
części wstępne i metodyczne artykułów w katalogu artykuły/.
Tekst z pracy magisterskiej i artykułów wymaga pogłębienia, ujednolicenia notacji z config/symbols.tex oraz weryfikacji źródeł.
Rozprawa Michała Kordasza w katalogu ,,doktorat Michał'' jest wyłącznie wzorcem układu i poziomu szczegółowości; nie wolno kopiować z niej tekstu, wzorów ani rysunków.
}

% !TeX root = ./../main.tex

\section{Podstawy informatyki kwantowej}\label{sec:tb:quantum}
\agentinfo{
Cel: wprowadzić minimalny, ale kompletny aparat potrzebny do zrozumienia metod z rozdziałów 4--8.
Poziom: czytelnik jest informatykiem bez przygotowania z mechaniki kwantowej.
Źródła własne: \agentref{artykuły/magisterka_werner_final/chapters/01.tex} (sekcje: Notacja Diraca, Nowe sformułowanie mechaniki kwantowej, Podstawy informatyki kwantowej).
Źródła podręcznikowe: Nielsen i Chuang oraz pozycje z bibliography/biblio.bib; brakujące pozycje dodać po weryfikacji.
Rysunki (sfera Blocha, obwody) wykonać samodzielnie w TikZ i quantikz2, w katalogu figures/chp_theoretical_background.
}

\subsection{Kubit, rejestry kwantowe i notacja Diraca}\label{subsec:tb:qubit}
\agentinfo{
Bit a kubit, superpozycja, przestrzeń Hilberta, notacja Diraca, iloczyn tensorowy, rejestry wielokubitowe, splątanie, sfera Blocha.
}

\subsection{Bramki i obwody kwantowe}\label{subsec:tb:gates}
\agentinfo{
Ewolucja unitarna, bramki jednokubitowe (Pauliego, Hadamarda, obroty), bramki przesunięcia fazy, bramki kontrolowane i wielokrotnie kontrolowane, obwód kwantowy, głębokość obwodu.
Techniki używane dalej w pracy: phase kickback, uncomputation, definiowanie operatora macierzą i jego transpilacja do bramek natywnych urządzenia.
Zwrócić uwagę, że transpilacja zmienia liczbę i rodzaj bramek, co wpływa na błąd; do tego odwołują się rozdziały 4 i 5.
}

\subsection{Faza globalna i lokalna, interferencja}\label{subsec:tb:phase}
\agentinfo{
Rozróżnienie fazy globalnej (nieobserwowalnej) i lokalnej (względnej).
Interferencja jako mechanizm zamiany informacji fazowej na prawdopodobieństwa pomiaru.
To pojęcie jest fundamentem kodowania LPIQE i metody PDU, dlatego wymaga osobnego, starannego wprowadzenia z przykładem jednokubitowym.
}

\subsection{Pomiar i estymacja rozkładu wyników}\label{subsec:tb:measurement}
\agentinfo{
Postulat pomiaru, pomiar w bazie obliczeniowej, wartość oczekiwana.
Wielokrotne wykonanie obwodu (shots), histogram wyników jako estymator rozkładu prawdopodobieństwa.
Błąd statystyczny wynikający ze skończonej liczby powtórzeń i jego odróżnienie od błędu sprzętowego; to rozróżnienie jest potrzebne przy interpretacji wyników w rozdziałach 4--6 i 8.
}

\subsection{Stany mieszane i macierz gęstości}\label{subsec:tb:density}
\agentinfo{
Stany czyste i mieszane, macierz gęstości, ślad częściowy, wierność stanów.
Wprowadzenie potrzebne do opisu szumu w kolejnym podrozdziale oraz do macierzy gęstości prezentowanych w rozdziale 4.
}

\subsection{Bramki kwantowe jako elementy grup Liego}\label{subsec:tb:lie}
\agentinfo{
Grupy SU(2), SO(3), SU(3) i ich algebry Liego, generatory, mapa wykładnicza.
Bramka obrotu jako eksponenta generatora: U(theta) = exp(-i theta G / 2).
Jest to podstawa rozdziału 8 (zaburzenia generatorów, sieci oparte na grupach Liego).
Źródła: \agentref{artykuły/quantum_noise_elsevier_submission/main.tex} (sekcja Theoretical Foundations) \agentref{artykuły/noisy_lie_group_learners/main.tex} (sekcja Lie Groups, Lie Algebras, and Neural Network Layers).
Zakres ograniczyć do tego, co faktycznie jest używane w rozdziale 8.
}

% !TeX root = ./../main.tex

\section{Szum i błędy w komputerach kwantowych}\label{sec:tb:noise}
\agentinfo{
Stan (PHD-TS-95, 2026-11-27): podrozdział skrócony i poprawiony po recenzji \agentref{eksperymenty/recenzje/rozdzial-2-sekcje-2.1-2.2.md}; lista cięć (w tym rysunek z kulą Blocha, równanie Lindblada, miary QV i EPLG, tomografia zbioru bramek, wzór zaniku w teście RB) w \agentref{eksperymenty/recenzje/rozdzial-2-skracanie.md}.
Tabela kanałów tab:tb:noise:channels jest przeniesiona do dodatku \agentref{chapters/appendix_results_tb.tex} (sekcja app:results:tb); wiersz input tego pliku w appendix_results.tex dopisuje orkiestrator.
Etykiety używane przez inne pliki (nie usuwać): etykiety podsekcji, eq:tb:noise:affine i tab:tb:noise:channels (rozdział 8), eq:tb:noise:shrink (rozdział 7), eq:tb:noise:readout (podrozdział 2.3, dodatek A), eq:tb:noise:fringe (rozdziały 3 i 4), eq:tb:noise:ptm-blocks (plik archiwalny noise_structure.tex).
Notacja: parametry kanałów to q_d, q_a, q_phi (nie gamma, bo gamma to funkcja referencyjna), kąt błędu koherentnego to vartheta (nie epsilon, bo epsilon to funkcja PDU), rozkład na zaszumionym urządzeniu to makro probnoisy, macierz odczytu to makro confmat.
Kanał depolaryzujący ma dwie parametryzacje: lambda (Qiskit Aer) i q_d (PennyLane); dla jednego kubitu lambda = 4 q_d / 3.
Model warstwowy: szum po każdym bloku ansatzu, tak jak w rozdziałach 7 i 8 (słowo warstwa pochodzi z kodu).
Macierz transferu Pauliego, jej postać blokowa i liczba obwodów tomografii procesu służą tylko akapitowi par:lr:noise:generators (sekcja 3.5); nazwę wiersz afiniczny usunięto, bo sekcja 3.5 jej nie używa (ustalenie B1 recenzji dawnego rozdziału 9 dotyczy już tylko sekcji 3.5).
Liczby dla ibm_fez pochodzą z migawki kalibracyjnej z 2025-02-26 dołączonej do pakietu qiskit-ibm-runtime 0.49.0, a nie z dni eksperymentów; wartości z dni eksperymentów należą do dodatku \agentref{chapters/appendix_experiments.tex}.
Twierdzenie o układach fotonicznych sprawdzono tylko na poziomie streszczeń.
}

Rzeczywisty procesor kwantowy nie jest układem izolowanym o~dokładnych bramkach i~bezbłędnym pomiarze, jak założono w~podrozdziale~\ref{sec:tb:quantum}.
W~pracy \emph{szumem} nazywa się procesy fizyczne i~niedoskonałości sterowania zaburzające stan rejestru, a~\emph{błędem} ich skutek, czyli różnicę między operacją lub wynikiem otrzymanym a~idealnym.

\subsection{Źródła błędów}\label{subsec:tb:noise:sources}

Podstawowym źródłem błędów jest \gls{decoherence}, czyli utrata kwantowych własności układu wskutek niekontrolowanego oddziaływania z~otoczeniem~\cite{zurek_decoherence_2003,bib:quantum-decoherence}.
W~kubitach nadprzewodzących opisują ją czas relaksacji $\Tone$, z~którym zanika obsadzenie stanu $\ket{1}$, i~czas koherencji $\Ttwo$, z~którym zanikają koherencje (podrozdział~\ref{subsec:tb:density}), a~wraz z~nimi dane zapisane w~fazie~\cite{krantz_quantum_2019}; fazę niszczy zarówno relaksacja, jak i~czyste defazowanie o~stałej czasowej $\Tphi$, stąd $1/\Ttwo=1/(2\Tone)+1/\Tphi$.
Pozostałe źródła to błędy sterowania~\cite{krantz_quantum_2019}, błędy przygotowania stanu i~pomiaru (SPAM), \gls{crosstalk}, czyli zaburzanie kubitu przez operacje wykonywane w~jego sąsiedztwie~\cite{sarovar_detecting_2020,mundada2019crosstalk}, oraz \glslink{qubit-leakage}{wyciek}, czyli przejście układu do stanu spoza przestrzeni obliczeniowej~\cite{Wood2015}; wskutek dwóch ostatnich błąd obwodu nie jest prostym złożeniem błędów bramek, lecz zależy od wyboru kubitów fizycznych i~od kontekstu, w~którym bramka jest wykonywana.

Skalę tych zjawisk ilustruje migawka kalibracyjna 156-kubitowego procesora ibm\_fez z~26~lutego 2025~r., dołączona do pakietu qiskit-ibm-runtime w~wersji 0.49.0~\cite{qiskit_ibm_runtime_2026}: mediany wynoszą $\Tone=145~\mu\mathrm{s}$, $\Ttwo=88~\mu\mathrm{s}$, błąd bramki jednokubitowej $2{,}3\cdot10^{-4}$, błąd bramki dwukubitowej CZ $3{,}8\cdot10^{-3}$ (bez par oznaczonych jako niesprawne) i~błąd odczytu $7{,}6\cdot10^{-3}$, więc o~jakości wyniku decyduje przede wszystkim liczba bramek dwukubitowych po transpilacji (podrozdział~\ref{subsec:tb:gates}).

Parametry szumu zmieniają się w~czasie~\cite{burnett_decoherence_2019}, podczas gdy standardowe metody charakteryzacji zakładają stałość błędów~\cite{proctor_detecting_2020}.
Według dokumentacji IBM Quantum (stan na 26~listopada 2026~r.) poszczególne parametry kalibruje się między zadaniami użytkowników, gdy wykonywane mniej więcej co godzinę zadania monitorujące wykryją odchylenie, a~raz na dobę wykonuje się pomiary wzorcowe, z~których pochodzą publikowane wartości parametrów~\cite{ibm_calibration_docs_2026}; w~latach 2021--2022, w~okresie eksperymentów z~rozdziału~\ref{chapter:pdu}, deklarowane wartości błędów bramek zmieniały się według pracy~\cite{werner_experiment_driven_2022} raz na dobę (obserwacje innych autorów: podrozdział~\ref{sec:lr:noise}).
Metoda PDU zakłada, że między kolejnymi kalibracjami błąd urządzenia jest w~przybliżeniu ustalony (założenie sprawdzane w~podrozdziale~\ref{subsec:pdu:results:stability}), a~każda charakterystyka szumu, także korekta wyznaczona eksperymentalnie, dotyczy konkretnych kubitów fizycznych i~ma ograniczony czas ważności.

\paragraph{Era NISQ i~platformy sprzętowe.}
Obecne procesory należą do okresu określanego jako \gls{nisq}~\cite{preskill_quantum_2018}: szum bramek ogranicza rozmiar obwodu, który można wykonać wiarygodnie, a~liczba kubitów nie wystarcza do pełnej korekcji błędów.
Eksperymenty w~tej pracy wykonano na procesorach IBM Quantum z~kubitami nadprzewodzącymi~\cite{krantz_quantum_2019}, w~których połączenia tylko między sąsiednimi kubitami wymuszają dodatkowe bramki dwukubitowe przy transpilacji; parametry użytych urządzeń zebrano w~dodatku~\ref{appendix:experiments}.
Kodowanie LPIQE (podrozdział~\ref{subsec:tb:qimg:lpiqe}) zaprojektowano dla układów fotonicznych, w~których fazę zadaje przesuwnik fazy, a~głównymi źródłami błędów są straty i~niepełna nierozróżnialność fotonów~\cite{slussarenko_photonic_2019,taballione_20-mode_2022}, więc wyników z~urządzeń nadprzewodzących nie można wprost przenosić na tę platformę.

\subsection{Formalizm kanałów kwantowych}\label{subsec:tb:noise:channels}

Zmiany stanu kubitu oddziałującego z~otoczeniem nie opisuje macierz unitarna, lecz \gls{quantum-channel}: liniowe odwzorowanie $\channel$ macierzy gęstości, całkowicie dodatnie i~zachowujące ślad (CPTP)~\cite{NielsenChuang2010,bib:quantum_noise}.
Każdy kanał ma postać sumy operatorowej $\channel(\rho)=\sum_{k}\kraus_k\,\rho\,\kraus_k^{\dagger}$ z~operatorami Krausa $\kraus_k$ spełniającymi $\sum_{k}\kraus_k^{\dagger}\kraus_k=\identity$.
Bramka idealna $U$ jest kanałem o~jednym operatorze Krausa, $\mathcal{U}(\rho)=U\rho U^{\dagger}$, a~bramkę zaszumioną przedstawia się jako złożenie $\errch\circ\mathcal{U}$ bramki idealnej i~\emph{kanału błędu} $\errch$.

\paragraph{Błędy koherentne i~niekoherentne.}
\Gls{coherent-error} jest unitarnym kanałem błędu, na przykład obrotem o~dodatkowy kąt $\cohangle$ wynikającym z~niedokładnej kalibracji impulsu; jest jednakowy w~każdym wykonaniu obwodu, nie zamienia stanu czystego w~mieszany i~daje się odwrócić, o~ile jest znany.
\Gls{incoherent-error}, nazywany też stochastycznym, jest kanałem, którego nie da się zapisać jednym operatorem Krausa, na przykład kanałem Pauliego, w~którym z~prawdopodobieństwem $q_j$ na stan działa macierz $\sigma_j\in\{\identity,\sigma_x,\sigma_y,\sigma_z\}$; na ogół zamienia on stany czyste w~mieszane, a~jego odwrotność nie jest kanałem.
W~błędach koherentnych sumują się amplitudy, więc $k$ powtórzeń bramki z~nadmiarowym obrotem $\cohangle$ daje prawdopodobieństwo błędnego wyniku rzędu $k^2\cohangle^2/4$, podczas gdy w~błędach niekoherentnych sumują się prawdopodobieństwa i~wzrost jest liniowy względem $k$~\cite{Wallman2015,Hashim2023}.

\paragraph{Kanały unitalne i~nieunitalne.}
Kanał jest \glslink{unital-channel}{unitalny}, jeżeli nie zmienia stanu maksymalnie mieszanego, $\channel(\identity/d)=\identity/d$, gdzie $d=2^n$; kanały unitarne i~kanały Pauliego są unitalne, a~nieunitalna jest na przykład relaksacja, która w~granicy długiego czasu każdy stan sprowadza do $\ket{0}$.
Dla jednego kubitu, w~reprezentacji~\eqref{eq:tb01:blochrho}, każdy kanał jest przekształceniem afinicznym wektora Blocha,
\begin{equation}\label{eq:tb:noise:affine}
\blochvec\;\longmapsto\;\ptmblock\,\blochvec+\ptmshift,
\end{equation}
z~rzeczywistą macierzą $\ptmblock$ o~wymiarach $3\times3$ i~wektorem przesunięcia $\ptmshift$~\cite{NielsenChuang2010}, a~unitalność oznacza $\ptmshift=\mathbf{0}$.
Dla bramki unitarnej $\ptmblock$ jest macierzą obrotu, szum niekoherentny zmniejsza jej wartości osobliwe poniżej jedności, czyli ściąga kulę Blocha, a~kanał nieunitalny przesuwa ponadto jej środek.

\subsection{Podstawowe modele szumu}\label{subsec:tb:noise:models}

Modele szumu używane w~tej pracy buduje się z~kanałów jednokubitowych~\cite{NielsenChuang2010}, których operatory Krausa oraz macierze $\ptmblock$ i~wektory $\ptmshift$ zestawia tabela~\ref{tab:tb:noise:channels} w~dodatku~\ref{app:results:tb}.
Tłumienie amplitudy o~parametrze $\qamp$ modeluje relaksację i~jest jedynym z~nich kanałem nieunitalnym, a~tłumienie fazy o~parametrze $\qphase$, które mnoży składowe $r_x$ i~$r_y$ wektora Blocha przez $c_{\varphi}=\sqrt{1-\qphase}$, modeluje czyste defazowanie; odwrócenie fazy, czyli kanał Pauliego z~macierzą $\sigma_z$ i~prawdopodobieństwem $q$, jest tym samym kanałem w~innej parametryzacji, $1-2q=c_{\varphi}$.
Złożenie obu kanałów, nazywane w~Qiskit Aer relaksacją termiczną (tu w~temperaturze zerowej), mnoży po czasie $t$ składowe $r_x$ i~$r_y$ przez $\mathrm{e}^{-t/\Ttwo}$, a~najprostszym modelem błędu koherentnego jest odstrojenie, czyli stały obrót $R_z(\cohangle)$.

Kanał depolaryzujący o~parametrze $\depolparam$ zastępuje stan $n$ kubitów mieszaniną $(1-\depolparam)\,\rho+\depolparam\,\identity/2^n$, a~dla jednego kubitu jest kanałem Pauliego, w~którym każda z~macierzy $\sigma_x$, $\sigma_y$, $\sigma_z$ działa z~prawdopodobieństwem $\qdep/3$, przy czym $\depolparam=\tfrac{4}{3}\qdep$; parametru $\depolparam$ używa Qiskit Aer~\cite{qiskit_aer_noise_docs_2026}, a~$\qdep$ PennyLane (operacja DepolarizingChannel)~\cite{bergholm_pennylane_2018}, i~nie wolno ich utożsamiać.
Po działaniu tego kanału na cały rejestr rozkład wyników jest mieszaniną rozkładu idealnego i~jednostajnego,
\begin{equation}\label{eq:tb:noise:shrink}
\probnoisy(\ket{x})=\distshrink\,\prob_{\mathrm{id}}(\ket{x})+\frac{1-\distshrink}{2^n},\qquad \distshrink=1-\depolparam,
\end{equation}
czyli odchylenia od rozkładu jednostajnego maleją w~stosunku $\distshrink$, nazywanym dalej \glslink{distribution-shrinkage}{współczynnikiem skurczenia rozkładu}.
Błąd odczytu opisuje się klasycznie, jako losową zamianę etykiet wyników~\cite{maciejewski_mitigation_2020}:
\begin{equation}\label{eq:tb:noise:readout}
\probnoisy(\ket{y})=\sum_{x\in\{0,1\}^n}\confmat_{yx}\,\prob(\ket{x}),\qquad \confmat_{yx}=q_{y|x},
\end{equation}
gdzie $\prob$ jest rozkładem przed odczytem, a~$q_{y|x}$ prawdopodobieństwem odczytania wyniku $y$ dla rejestru w~stanie $\ket{x}$; wyznaczanie i~odwracanie macierzy $\confmat$ opisano w~podrozdziale~\ref{subsec:tb:qec:mitigation}.

\paragraph{Model warstwowy.}
W~rozdziałach~\ref{chapter:learned_encoding} i~\ref{chapter:classification} używany jest fenomenologiczny model, w~którym po każdym bloku obwodu, czyli powtarzanym fragmencie ansatzu (w~kodzie: warstwie, stąd nazwa), na każdy kubit działają odstrojenie $R_z(\cohangle)$, tłumienie fazy o~parametrze $\qphase$ i~kanał depolaryzujący o~parametrze $\qdep$.
Symulacje z~tym modelem nie zawierają błędu odczytu~\eqref{eq:tb:noise:readout} (zawiera go tylko model budowany z~danych kalibracyjnych, opisany niżej), a~w~części konfiguracji wyłączone jest tłumienie fazy albo odstrojenie.
Te trzy odwzorowania są przemienne, więc ich kolejność w~obrębie bloku nie wpływa na wynik.
W~eksperymencie interferencyjnym z~przykładu~\ref{ex:tb01:interference} jedna dawka takiego szumu przed drugą bramką Hadamarda daje
\begin{equation}\label{eq:tb:noise:fringe}
\probnoisy(\ket{0})=\tfrac{1}{2}\bigl[1+v\cos(\varphi+\cohangle)\bigr],\qquad v=c_{\varphi}\bigl(1-\tfrac{4}{3}\qdep\bigr),
\end{equation}
czyli składowa koherentna przesuwa argument funkcji, a~składowe niekoherentne zmniejszają widoczność interferencji $v$ (podrozdział~\ref{subsec:tb:density}).

\paragraph{Szum a~składowe błędu wyniku.}
W~rozkładzie~\eqref{eq:tb01:errordecomp} błąd sprzętowy $\probnoisy-\prob_{\mathrm{id}}$ jest w~całości wyznaczony przez kanał obwodu i~macierz odczytu, a~wchodzą do niego zarówno błędy koherentne, jak i~niekoherentne: w~równaniu~\eqref{eq:tb:noise:fringe} przesunięcie $\cohangle$ i~widoczność $v$ łącznie określają deterministyczną funkcję fazy, powtarzalną, dopóki parametry urządzenia pozostają ustalone.
Podział na błędy koherentne i~niekoherentne nie pokrywa się zatem z~podziałem na składową systematyczną i~statystyczną, która wynika wyłącznie ze skończonej liczby pomiarów.
Rozróżnienie wyznacza granice klasycznej korekty wyniku (rozdziały~\ref{chapter:pdu}, \ref{chapter:learned_encoding} i~\ref{chapter:classification}): korekta addytywna, czyli odjęcie oszacowania składowej systematycznej, nie zmienia wariancji składowej statystycznej poza niepewnością samego oszacowania, natomiast korekta odwracająca skurczenie~\eqref{eq:tb:noise:shrink} wymaga dzielenia przez $\distshrink<1$, a~zatem zwiększa odchylenie standardowe składowej statystycznej $1/\distshrink$ razy.

\paragraph{Modele szumu w~symulatorach.}
W~modelu budowanym w~Qiskit Aer z~danych kalibracyjnych urządzenia po każdej bramce działa kanał depolaryzujący, a~po nim relaksacja termiczna o~parametrach $\Tone$, $\Ttwo$ i~czasie trwania bramki, przy czym parametr depolaryzacji dobiera się tak, aby średni błąd złożenia był równy błędowi bramki z~kalibracji; odczyt opisują niezależne jednokubitowe macierze $\confmat$~\cite{qiskit_aer_noise_docs_2026}.
Model ten nie zawiera błędów koherentnych, przesłuchu, wycieku, korelacji błędów odczytu ani dryfu, a~parametry dopasowane do wyników urządzenia różnią się od kalibracyjnych~\cite{georgopoulos_modeling_2021}, dlatego symulacja z~szumem jest przybliżeniem, a~nie zamiennikiem pomiaru na sprzęcie.

\subsection{Charakteryzacja szumu}\label{subsec:tb:noise:characterisation}

Tomografia procesu kwantowego (QPT) odtwarza kanał z~tomografii stanów wyjściowych, czyli z~pomiarów w~kilku bazach, dla zbioru stanów wejściowych rozpinającego przestrzeń macierzy gęstości~\cite{bib:quantum_tomography,ChuangNielsen1997,Poyatos1997}.
Dla jednego kubitu wystarczają cztery stany wejściowe i~trzy bazy pomiarowe, czyli dwanaście obwodów, lecz liczba parametrów kanału rośnie jak $16^n$~\cite{NielsenChuang2010,MohseniQPTReview}, a~badanej operacji przypisuje się także błędy SPAM~\cite{Merkel2013}.
Wynik QPT przedstawia się jako \gls{ptm} (PTM)~\cite{greenbaum2015introduction}, która dla jednego kubitu działa na wektor $(1,r_x,r_y,r_z)^{\top}$ i~ma postać blokową:
\begin{equation}\label{eq:tb:noise:ptm-blocks}
R_{ab}=\frac{1}{2}\Tr\bigl[\sigma_a\,\channel(\sigma_b)\bigr],\quad a,b\in\{0,x,y,z\},\quad \sigma_0=\identity,
\qquad
\ptm=\begin{pmatrix} R_{00} & \ptmrow^{\top}\\ \ptmshift & \ptmblock \end{pmatrix},
\end{equation}
gdzie $\ptmblock$ i~$\ptmshift$ pochodzą z~równania~\eqref{eq:tb:noise:affine}.
Macierz $\ptmblock$ nazywa się \emph{blokiem unitalnym} (dla bramki $\mathrm{X}$ jest równa $\operatorname{diag}(1,-1,-1)$), a~wektor $\ptmshift$ częścią nieunitalną kanału~\cite{greenbaum2015introduction,Wallman2015}.
Dla kanału zachowującego ślad $R_{00}=1$ i~$\ptmrow=\mathbf{0}$; niezerowe elementy $\ptmrow$ są możliwe tylko w~rekonstrukcji bez narzuconego zachowania śladu~\cite{Wallman2015}.

Błędy bramek w~danych kalibracyjnych wyznacza się testem z~losowymi sekwencjami bramek (ang. \emph{randomized benchmarking}, RB)~\cite{Magesan2011,magesan2012characterizing,ibm_calibration_docs_2026}, który przy błędzie niezależnym od bramki uśrednia kanał błędu do kanału depolaryzującego o~tej samej średniej wierności; uzasadnia to użycie tego kanału w~modelach symulatorów, ale sprawia też, że wynik RB nie rozróżnia błędów koherentnych i~niekoherentnych~\cite{Wallman2015,ProctorWhat2017}.

% !TeX root = ./../main.tex

\section{Korekcja, mitygacja i redukcja błędów}\label{sec:tb:qec}
\agentinfo{
Uzgodnienia (PHD-TS-14, 2026-11-26): tutaj definicje i mechanizmy metod; stan badań i wyniki prac w \agentref{chapters/literature_review.tex}.
Terminologia całej pracy: korekcja błędów oznacza wyłącznie QEC, poprawianie wyniku po stronie klasycznej to korekta wyniku, mitygacja bez dopowiedzenia to metody QEM inne niż redukcja.
Metodę PDU definiuje rozdział 4; kryteria K1--K6 i tabela tab:tb03:cost służą dyskusjom w rozdziałach 4--6 i 8.
Notacja: skala szumu c, quasi-prawdopodobieństwa omega, koszt próbkowania makro pecnorm (gamma, kappa, lambda, eta i q są zajęte).
Zaadaptowane źródła własne autora: \agentref{artykuły/magisterka_werner_final/chapters/01.tex} \agentref{artykuły/Quantum_Errors_ICCS/main.tex} \agentref{artykuły/CoDIT_2025_final_paper/root.tex} \agentref{artykuły/Quantum_journal_GAN__Copy_/article.tex}.
}

Błąd sprzętowy, czyli systematyczna składowa błędu wyniku w rozkładzie~\eqref{eq:tb01:errordecomp}, nie maleje ze wzrostem liczby powtórzeń obwodu, więc trzeba mu przeciwdziałać innymi środkami.
W tym podrozdziale zdefiniowano trzy rodziny metod, które temu służą, opisano mechanizmy metod bazowych i wprowadzono kryteria kosztu stosowane w dalszych rozdziałach.
Stan badań i wyniki eksperymentalne omówiono w rozdziale~\ref{chapter:literature_review}.

\subsection{Terminologia przyjęta w pracy}\label{subsec:tb:qec:taxonomy}

W pracy rozróżnia się korekcję, mitygację i redukcję błędów; o przynależności metody decyduje to, czy zmienia ona obwód kwantowy zadania i jakich dodatkowych zasobów wymaga (rys.~\ref{fig:tb03:taxonomy}).

\begin{Definition}[Kwantowa korekcja błędów]\label{def:tb03:qec}
\glsdisp{error-correction}{Kwantową korekcją błędów} (ang. \emph{quantum error correction}, QEC) nazywa się metody, w których stan $k$ \glsdisp{logical-qubit}{kubitów logicznych} jest zakodowany nadmiarowo w stanie $n>k$ kubitów fizycznych, a błędy są wykrywane w trakcie obliczeń pomiarem \glsdisp{syndrome}{syndromu} i odwracane operacją zależną od jego wyniku, bez pomiaru zakodowanej informacji.
\end{Definition}

\begin{Definition}[Mitygacja błędów kwantowych]\label{def:tb03:qem}
\glsdisp{error-mitigation}{Mitygacją błędów kwantowych} (ang. \emph{quantum error mitigation}, QEM) nazywa się metody, które zmniejszają wywołane szumem obciążenie estymatora wartości oczekiwanej lub rozkładu wyników przez klasyczną obróbkę wyników wielu wykonań obwodu, jego zmodyfikowanych wariantów lub obwodów kalibracyjnych, bez kodowania informacji w kubitach logicznych.
\end{Definition}

\begin{Definition}[Redukcja błędów kwantowych]\label{def:tb03:reduction}
\glsdisp{error-reduction}{Redukcją błędów kwantowych} nazywa się w niniejszej pracy metody, które korygują wynik obliczeń kwantowych wyłącznie po stronie klasycznej, modelem korekty wyznaczonym eksperymentalnie dla danego zadania obliczeniowego i urządzenia, przy czym obwód zadania jest wykonywany bez zmian: bez dodatkowych kubitów i bramek, a dodatkowe wykonania służą tylko kalibracji modelu.
\end{Definition}

\begin{figure}[tb]
  \centering
  % !TeX root = ./../../main.tex
% Miejsce ingerencji trzech rodzin metod w proces obliczeń (podrozdział 2.3, rys. fig:tb03:taxonomy).
% Rysunek własny (TikZ): górny wiersz to etapy obliczeń opartych na próbkowaniu,
% trzy pasy poniżej leżą pod etapami, które zmienia korekcja, mitygacja i redukcja błędów.
\begin{tikzpicture}[
  font=\footnotesize,
  stage/.style={draw, rounded corners=2pt, minimum height=0.8cm, text width=3.0cm, align=center, inner sep=2pt},
  band/.style={draw, fill=black!8, inner sep=2.5pt, align=center, minimum height=0.8cm},
  >={Stealth[length=2mm]}]
  % etapy obliczeń
  \node[stage] (circ) at (0,0)    {obwód kwantowy\\zadania};
  \node[stage] (meas) at (3.7,0)  {wykonania obwodu\\i~pomiar};
  \node[stage] (post) at (7.4,0)  {klasyczna obróbka\\statystyk};
  \node[stage] (res)  at (11.1,0) {wynik zadania};
  \draw[->] (circ) -- (meas);
  \draw[->] (meas) -- (post);
  \draw[->] (post) -- (res);
  % pasy: zakres ingerencji metod (szerokość pasa = etapy, które metoda zmienia)
  \node[band, text width=6.6cm] (qec) at (1.85,-1.05)
    {\textbf{korekcja (QEC):} nadmiarowe kubity,\\pomiar syndromu, operacje naprawcze};
  \node[band, text width=10.3cm] (qem) at (3.7,-2.0)
    {\textbf{mitygacja (QEM):} warianty obwodu lub obwody kalibracyjne,\\dodatkowe wykonania, obróbka zebranych statystyk};
  \node[band, text width=6.6cm] (red) at (9.25,-2.95)
    {\textbf{redukcja:} korekta wyniku modelem\\z~kalibracji; obwód zadania bez zmian};
\end{tikzpicture}

  \caption[Miejsce ingerencji korekcji, mitygacji i redukcji błędów w proces obliczeń]{Miejsce ingerencji trzech rodzin metod w proces obliczeń; pas każdej rodziny leży pod etapami, które ta rodzina zmienia.}
  \label{fig:tb03:taxonomy}
\end{figure}

Korekcja usuwa błędy ze stanu kwantowego w każdym wykonaniu obwodu.
Mitygacja i redukcja tego nie robią: w pojedynczym wykonaniu poziom szumu pozostaje taki sam albo wyższy, a poprawia się dopiero wynik wyznaczony ze zbioru wykonań~\cite{cai2023quantum}.
Tytułowe ,,ograniczanie błędów'' obejmuje właśnie te dwie rodziny.
Do redukcji należą metody będące przedmiotem pracy: funkcja odwracania zniekształceń fazy (PDU), w której model korekty wyznacza się z pomiarów funkcji referencyjnej (rozdział~\ref{chapter:pdu}), oraz korekta odtworzonych obrazów modelem generatywnym uczonym na parach obraz zaszumiony--obraz wzorcowy (rozdział~\ref{chapter:dl_error_reduction}).
Termin korekcja błędów jest w pracy zastrzeżony dla QEC, a poprawianie wyniku po stronie klasycznej, będące krokiem metod redukcji, nazywane jest korektą wyniku.

Granice między rodzinami są częściowo umowne.
Mitygacja nie ma powszechnie przyjętej definicji; definicja robocza z przeglądu~\cite{cai2023quantum}, za którą idzie definicja~\ref{def:tb03:qem}, obejmuje każdą obróbkę wyników zbioru wykonań obwodu zmniejszającą obciążenie wywołane szumem.
W tym szerokim sensie redukcja jest szczególnym przypadkiem mitygacji, a angielski termin \emph{error reduction} nie oznacza w literaturze odrębnej klasy metod.
W pracy wyodrębniono ją, ponieważ nie modyfikuje obwodu zadania i jest zorientowana na problem: korektę wyznacza się dla wyniku całego zadania, bez modelu szumu bramek lub pomiaru.
Słowo mitygacja bez dopowiedzenia oznacza dalej metody QEM inne niż redukcja.
Blisko redukcji leży mitygacja błędu odczytu (podrozdział~\ref{subsec:tb:qec:mitigation}), która także nie zmienia obwodu zadania i sprowadza się do kalibracji oraz obliczeń klasycznych~\cite{maciejewski_mitigation_2020,bravyi_mitigating_2021}.
Kalibruje się w niej jednak sam pomiar, niezależnie od obwodu, więc korekta przenosi się na dowolne zadanie, ale obejmuje tylko błąd odczytu; w redukcji kalibruje się wynik całego zadania, więc korekta obejmuje łącznie skutki błędów bramek, dekoherencji i odczytu, lecz nie przenosi się na inne zadania.
Jeszcze bliżej leżą metody, w których korektę wyniku kalibruje się na obwodach o znanym wyniku~\cite{jattana_general_2020,czarnik_error_2021,lolur_reference-state_2023}; ich zestawienie z metodą PDU zawiera podrozdział~\ref{sec:lr:qem} (tabela~\ref{tab:lr:qem:calibrated}).
\agentinfo{
Dopisek orkiestratora (2026-11-27): zdanie o metodach kalibrowanych na obwodach o znanym wyniku dodano po zebraniu literatury do sekcji 3.3, żeby podrozdział nie był sprzeczny z przeglądem literatury.
Decyzja autora z 2026-11-27: PDU jest metodą rozwijaną w ramach tego doktoratu i ma być opisana jako ważny wkład rozprawy, z wyróżnikami podanymi w sekcji 3.3 (korekta indeksowana znanym argumentem i zbadanie stabilności błędu między kalibracjami).
Nadal do decyzji autora: czy redukcję opisywać jako kategorię obejmującą te metody (definicję def:tb03:reduction spełniają także one) oraz czy dodać dla nich wiersz w tabeli tab:tb03:cost, która przypisuje całej mitygacji zorientowanie na bramkę.
}
Poza trzema definicjami pozostają metody pośrednie: weryfikacja symetrii, która stosuje sprawdzenia podobne do pomiaru syndromu, lecz wykonania z wykrytym błędem odrzuca, zamiast je naprawiać, oraz dynamiczne odsprzęganie, technika sterowania działająca bez pomiaru i obróbki wyników, nazywana tłumieniem błędów~\cite{cai2023quantum,ezzell_dynamical_2023}; omówiono je razem z mitygacją.
We wcześniejszych pracach autora metodę PDU nazywano także korekcją błędów~\cite{werner_experiment_driven_2022}, a jej działanie określano jako mitygację~\cite{werner_magisterka_2022}; w rozprawie stosowany jest wyłącznie termin redukcja.

\subsection{Kwantowa korekcja błędów}\label{subsec:tb:qec:codes}

Klasyczne kody korekcyjne powielają bity i rozstrzygają o ich wartości głosowaniem, czego nie da się przenieść wprost na kubity z trzech powodów~\cite{NielsenChuang2010,bib:error_corr_intr}.
Twierdzenie o zakazie klonowania wyklucza istnienie operacji unitarnej $U$ spełniającej $U\ket{\psi}\ket{0}=\ket{\psi}\ket{\psi}$ dla każdego stanu $\ket{\psi}$, więc nieznanego stanu nie można powielić~\cite{wootters_single_1982}.
Pomiar kubitów danych niszczy superpozycję (podrozdział~\ref{subsec:tb:measurement}), a błędy tworzą zbiór ciągły: obok odwrócenia bitu (bramka $\mathrm{X}$) występują odwrócenie fazy ($\mathrm{Z}$), ich złożenie ($\mathrm{Y}$) oraz dowolne kombinacje liniowe tych operatorów.

Kody kwantowe omijają te przeszkody: stan zapisuje się w splątanym stanie wielu kubitów zamiast w kopiach, a mierzy się wyłącznie syndrom, czyli wielkości zależne od błędu, lecz nie od zakodowanych amplitud.
Najprostszym przykładem jest trzykubitowy kod powtórzeniowy, w którym stan $\alpha\ket{0}+\beta\ket{1}$ zapisuje się jako $\alpha\ket{000}+\beta\ket{111}=\alpha\ket{0_L}+\beta\ket{1_L}$, gdzie $\ket{0_L}$ i $\ket{1_L}$ są stanami bazowymi kubitu logicznego.
Syndrom tworzą parzystości sąsiednich par kubitów, czyli wyniki pomiaru operatorów $\mathrm{Z}_1\mathrm{Z}_2$ i $\mathrm{Z}_2\mathrm{Z}_3$, odczytywane za pomocą dwóch kubitów pomocniczych (rys.~\ref{fig:tb03:repetition}).
Odwrócenie jednego bitu przenosi stan do jednej z trzech podprzestrzeni ortogonalnych do podprzestrzeni kodowej, a każdej z nich odpowiada inna wartość syndromu, więc błąd można rozpoznać i odwrócić bramką $\mathrm{X}$ bez poznania $\alpha$ i $\beta$.
Przy niezależnych odwróceniach bitów z prawdopodobieństwem $p$ kod zawodzi dopiero przy co najmniej dwóch błędach, czyli z prawdopodobieństwem $3p^2-2p^3$, mniejszym od $p$ dla $p<1/2$; przed odwróceniem fazy nie chroni.

\begin{figure}[tb]
  \centering
  % !TeX root = ./../../main.tex
% Trzykubitowy kod powtórzeniowy z pomiarem syndromu (podrozdział 2.3, rys. fig:tb03:repetition).
% Rysunek własny (quantikz2): kodowanie, kanał szumu, pomiar parzystości Z1Z2 i Z2Z3
% na dwóch kubitach pomocniczych, operacja naprawcza R sterowana klasycznie wynikiem syndromu.
\begin{quantikz}[row sep={0.55cm,between origins}, column sep=0.34cm, font=\small]
\lstick{$\ket{\psi}$} & \ctrl{1} & \ctrl{2} & \gate[3]{\channel} & \ctrl{3} &          &          &          &          & \gate[3]{R} & \rstick[3]{$\alpha\ket{000}+\beta\ket{111}$} \\
\lstick{$\ket{0}$}    & \targ{}  &          &                    &          & \ctrl{2} & \ctrl{3} &          &          &             & \\
\lstick{$\ket{0}$}    &          & \targ{}  &                    &          &          &          & \ctrl{2} &          &             & \\
\lstick{$\ket{0}$}    &          &          &                    & \targ{}  & \targ{}  &          &          & \meter{} & \ctrl[vertical wire=c]{-1}\setwiretype{c} & \rstick{$s_1$} \\
\lstick{$\ket{0}$}    &          &          &                    &          &          & \targ{}  & \targ{}  & \meter{} & \ctrl[vertical wire=c]{-1}\setwiretype{c} & \rstick{$s_2$}
\end{quantikz}

  \caption[Trzykubitowy kod powtórzeniowy z pomiarem syndromu]{Trzykubitowy kod powtórzeniowy: kodowanie stanu $\ket{\psi}=\alpha\ket{0}+\beta\ket{1}$, kanał szumu $\channel$, pomiar syndromu $(s_1,s_2)$ na kubitach pomocniczych i sterowana nim operacja naprawcza $R$.}
  \label{fig:tb03:repetition}
\end{figure}

Pomiar syndromu rzutuje dowolny błąd pojedynczego kubitu na jeden z operatorów $\identity$, $\mathrm{X}$, $\mathrm{Y}$, $\mathrm{Z}$, dlatego kod korygujący odwrócenia bitu i fazy koryguje każdy błąd jednokubitowy~\cite{NielsenChuang2010}.
Pierwszym takim kodem był dziewięciokubitowy kod Shora~\cite{bib:shor_code}, w którym każdy z trzech kubitów kodu powtórzeniowego chroniącego fazę jest zapisany trzema kubitami kodu chroniącego bit.
Najmniejszy kod o tej własności wykorzystuje pięć kubitów~\cite{bib:perfect_quantum_error_correction_code}: jeżeli brakowi błędu i każdemu z $3n$ błędów jednokubitowych ma odpowiadać osobna para ortogonalnych stanów, to musi zachodzić $2(3n+1)\le2^n$, skąd $n\ge5$.

Wymienione kody należą do rodziny kodów stabilizatorowych~\cite{gottesman_stabilizer_1997,NielsenChuang2010}.
Kod taki określa grupa stabilizatorów $\stabgroup$: przemienna, niezawierająca $-\identity$ podgrupa grupy Pauliego, czyli grupy iloczynów tensorowych operatorów $\identity$, $\mathrm{X}$, $\mathrm{Y}$, $\mathrm{Z}$ na $n$ kubitach z czynnikami $\pm1$, $\pm\mathrm{i}$.
Podprzestrzeń kodową tworzą stany spełniające $S\ket{\psi}=\ket{\psi}$ dla każdego $S\in\stabgroup$; przy $n-k$ niezależnych generatorach grupy mieści ona $k$ kubitów logicznych.
Syndromem jest ciąg wyników pomiaru generatorów: błąd będący operatorem Pauliego zmienia na $-1$ wyniki tych generatorów, z którymi antykomutuje.
Kod oznacza się symbolem $[[n,k,d]]$, w którym dystans $d$ jest najmniejszą liczbą kubitów, na które musi działać operator Pauliego, aby zmienić stan logiczny bez zmiany syndromu; taki kod, na przykład kod Shora $[[9,1,3]]$ albo pięciokubitowy $[[5,1,3]]$, koryguje dowolne błędy na $\lfloor(d-1)/2\rfloor$ kubitach.

W kodach powierzchniowych kubity leżą na dwuwymiarowej siatce, a generatorami są iloczyny operatorów $\mathrm{X}$ albo $\mathrm{Z}$ na co najwyżej czterech sąsiednich kubitach danych~\cite{kitaev_fault-tolerant_2003,bib:surface_codes_fowler}.
Każdy generator mierzy się kubitem pomocniczym sprzężonym tylko z najbliższymi sąsiadami, co sprzyja realizacji na urządzeniach o płaskiej topologii połączeń, a dystans rośnie z rozmiarem siatki.
Pomiary generatorów powtarza się cyklicznie w trakcie obliczeń, a klasyczny dekoder wyznacza z ciągu syndromów najbardziej prawdopodobny układ błędów.

O skalowalności korekcji rozstrzyga twierdzenie progowe: jeżeli prawdopodobieństwo błędu pojedynczej operacji $p$ jest mniejsze od wartości progowej $\pth$, to błąd obliczeń logicznych można uczynić dowolnie małym kosztem narzutu rosnącego polilogarytmicznie z rozmiarem obliczeń~\cite{aharonov_fault-tolerant_2008,knill_resilient_1998,NielsenChuang2010}.
Dla kodu powierzchniowego próg jest rzędu $1\%$ (w modelu błędów z pracy~\cite{bib:surface_codes_fowler} $\pth=0{,}57\%$), a poniżej progu prawdopodobieństwo błędu logicznego maleje wykładniczo z dystansem, $\plog\propto(p/\pth)^{\lfloor(d+1)/2\rfloor}$.
Ceną jest narzut zasobów: potrzebne są kubity pomiarowe, pomiary w trakcie obwodu i dekodowanie w czasie rzeczywistym, a według oszacowania z pracy~\cite{bib:surface_codes_fowler} kubit logiczny o użytecznie małym błędzie wymaga rzędu $10^3$--$10^4$ kubitów fizycznych.
Ograniczeniem jest też topologia: bramki działają tylko między sąsiadami na siatce, więc operacje na odległych kubitach logicznych wymagają dodatkowych kroków i dodatkowego miejsca~\cite{litinski_game_2019}.
Na urządzeniach ery NISQ (podrozdział~\ref{subsec:tb:noise:nisq}) kubity przeznaczone na korekcję pomniejszają i tak niewielką pulę dostępną dla obliczeń, co motywuje metody opisane dalej.

\subsection{Mitygacja błędów}\label{subsec:tb:qec:mitigation}

Metody mitygacji łączy wspólny schemat~\cite{cai2023quantum}: ze statystyk zaszumionych wykonań szacuje się wartość oczekiwaną obserwabli $O$ albo rozkład wyników obwodu idealnego.
Estymator po mitygacji ma mniejsze obciążenie, lecz większą wariancję niż estymator~\eqref{eq:tb01:estimator}, co trzeba skompensować większą liczbą wykonań obwodu.
Ten narzut próbkowania rośnie wykładniczo ze średnią liczbą błędów przypadających na jedno wykonanie obwodu, co ogranicza mitygację do obwodów o umiarkowanym rozmiarze.

\paragraph{Ekstrapolacja do zerowego szumu} (ang. \emph{zero-noise extrapolation}, ZNE)~\cite{richardson-extrapolation,quasi-probability}.
Ten sam obwód wykonuje się przy $m$ poziomach szumu $1=c_1<c_2<\dots<c_m$, gdzie $c=1$ oznacza poziom naturalny, i dla każdego wyznacza wartość $E(c_j)$ obserwabli.
Jeżeli zależność $E(c)$ przybliża wielomian stopnia $m-1$, to jego wartość w zerze, zwana ekstrapolacją Richardsona, wynosi
\begin{equation}\label{eq:tb03:zne}
  \hat{E}(0)=\sum_{j=1}^{m}\beta_j\,E(c_j),\qquad \beta_j=\prod_{l\neq j}\frac{c_l}{c_l-c_j}.
\end{equation}
Szum zwiększa się przez wydłużanie impulsów sterujących~\cite{kandala2021hardware} albo, bez dostępu do impulsów, przez zastępowanie bramki $G$ ciągiem $GG^{\dagger}G$~\cite{giurgica-tiron_digital_2020}.
Metoda zakłada, że szum daje się skalować w kontrolowany sposób, nie zmienia się w czasie zbierania danych i jest na tyle słaby, że $E(c)$ jest funkcją gładką.
Nie wymaga dodatkowych kubitów ani modelu szumu; kosztem jest $m$ serii wykonań, głębsze obwody przy powielaniu bramek oraz wzrost wariancji wyniku (wariancje pomiarów sumują się z wagami $\beta_j^2$).

\paragraph{Probabilistyczne znoszenie błędów} (ang. \emph{probabilistic error cancellation}, PEC), nazywane też metodą quasi-prawdopodobieństw~\cite{quasi-probability,endo_practical_2018}.
Idealną operację $\mathcal{U}$ przedstawia się jako kombinację liniową operacji $\mathcal{O}_i$ wykonalnych na zaszumionym urządzeniu:
\begin{equation}\label{eq:tb03:pec}
  \mathcal{U}=\sum_i\omega_i\,\mathcal{O}_i,\qquad \sum_i\omega_i=1,\qquad \pecnorm=\sum_i|\omega_i|\ge1,
\end{equation}
przy czym część współczynników $\omega_i$ jest ujemna, więc nie tworzą one rozkładu prawdopodobieństwa.
W każdym wykonaniu operację $\mathcal{O}_i$ losuje się z prawdopodobieństwem $|\omega_i|/\pecnorm$, a wynik pomiaru mnoży przez $\pecnorm\operatorname{sgn}(\omega_i)$; średnia tak ważonych wyników jest nieobciążonym estymatorem wartości idealnej.
Rozkład~\eqref{eq:tb03:pec} wymaga znajomości kanału szumu każdej bramki, wyznaczanego tomografią (podrozdział~\ref{subsec:tb:noise:characterisation}), i jego stałości w czasie.
W obwodzie złożonym z wielu bramek współczynniki $\pecnorm$ się mnożą, a liczba wykonań potrzebna do zachowania dokładności rośnie jak kwadrat ich iloczynu, czyli wykładniczo z liczbą bramek.

\paragraph{Mitygacja błędu odczytu macierzą kalibracyjną}~\cite{maciejewski_mitigation_2020,bravyi_mitigating_2021}.
Metoda odwraca klasyczny model błędu odczytu~\eqref{eq:tb:noise:readout}.
Kolumny macierzy $\confmat$ wyznacza się, przygotowując kolejne stany bazowe i mierząc je, a korektę uzyskuje się, działając macierzą $\confmat^{-1}$ na wektor zmierzonych częstości; jeżeli wynik zawiera wartości ujemne, zastępuje się go najbliższym rozkładem prawdopodobieństwa.
Pełna kalibracja rejestru $n$ kubitów wymaga $2^n$ obwodów, a przy założeniu niezależności błędów odczytu poszczególnych kubitów macierz ma tylko $2n$ parametrów.
Metoda zakłada, że błąd przygotowania stanów bazowych jest pomijalny, a macierz nie zmienia się między kalibracją a użyciem.
Obwód zadania pozostaje bez zmian, lecz korygowany jest wyłącznie błąd odczytu.

\paragraph{Dynamiczne odsprzęganie} (ang. \emph{dynamical decoupling}, DD)~\cite{viola_dynamical_1999,ezzell_dynamical_2023}.
W okresach, w których kubit czeka na kolejną bramkę, wstawia się ciąg szybkich impulsów, na przykład $\mathrm{X}$--$\mathrm{X}$ albo $\mathrm{X}$--$\mathrm{Y}$--$\mathrm{X}$--$\mathrm{Y}$, którego złożenie jest identycznością z dokładnością do fazy globalnej, lecz który uśrednia wolnozmienne sprzężenie kubitu z otoczeniem.
Wymaga to szumu zmieniającego się wolno w porównaniu z odstępem między impulsami oraz dokładnych impulsów; metoda dodaje bramki, ale nie wymaga dodatkowych kubitów, wykonań obwodu ani obróbki wyników.

Pozostałe metody mitygacji, w tym weryfikację symetrii i metody uczone na obwodach o znanym wyniku, omówiono w podrozdziale~\ref{sec:lr:qem}.

\subsection{Koszt zasobów jako kryterium porównania}\label{subsec:tb:qec:cost}

Opisane metody płacą za poprawę wyniku zasobami różnego rodzaju, dlatego w rozdziałach~\ref{chapter:pdu}--\ref{chapter:dl_error_reduction} oraz~\ref{chapter:classification} porównuje się je według sześciu kryteriów.
Są to: dodatkowe kubity, nieniosące danych zadania (K1); dodatkowe bramki, zwiększające głębokość obwodu zadania i jego podatność na szum (K2); dodatkowe wykonania obwodu, czyli obwody kalibracyjne i powtórzenia potrzebne do zachowania dokładności statystycznej (K3); koszt obliczeń klasycznych, w tym konieczność prowadzenia ich w czasie rzeczywistym (K4); założenia o szumie, czyli wymagany model szumu i jego stabilność między kalibracją a użyciem (K5; podrozdział~\ref{subsec:tb:noise:sources}); zorientowanie metody (K6).
Metoda zorientowana na bramkę przeciwdziała błędom elementów urządzenia (bramek, kubitów w~spoczynku, pomiaru) i stosuje się do dowolnego obwodu, a metoda zorientowana na problem wyznacza korektę dla wyniku konkretnego zadania, urządzenia i układu kubitów, traktując obwód jak czarną skrzynkę.
Kryteria K1 i K2 określają, jaka część zasobów urządzenia pozostaje dla właściwych obliczeń, co jest istotne w przetwarzaniu obrazów, które wymaga wielu kubitów i głębokich obwodów.
Zestawienie w tabeli~\ref{tab:tb03:cost} pokazuje, że redukcja błędów zajmuje położenie skrajne: poza kalibracją nie zużywa zasobów kwantowych, lecz ma najwęższy zakres stosowalności i opiera się na założeniu powtarzalności błędu, które trzeba sprawdzić eksperymentalnie (rozdział~\ref{chapter:pdu}).

\begin{table}[htb]
  \centering
  \caption[Koszt zasobów metod korekcji, mitygacji i redukcji błędów]{Koszt zasobów metod korekcji, mitygacji i redukcji błędów: dodatkowe kubity (K1), bramki (K2) i wykonania obwodu (K3), obliczenia klasyczne (K4), założenia o szumie (K5), zorientowanie na bramkę albo na problem (K6). Zestawienie na podstawie podrozdziałów~\ref{subsec:tb:qec:codes} i~\ref{subsec:tb:qec:mitigation}.}
  \label{tab:tb03:cost}
  \footnotesize
  \setlength{\tabcolsep}{3pt}
  \newcommand{\tbc}[1]{\raggedright\hspace{0pt}#1}
  \begin{tabular}{@{}p{2.5cm}p{0.6cm}p{2.3cm}p{2.5cm}p{2.3cm}p{2.95cm}p{1.45cm}@{}}
    \toprule
    \tbc{\textbf{Metoda}} & \tbc{\textbf{K1}} & \tbc{\textbf{K2}} & \tbc{\textbf{K3}} & \tbc{\textbf{K4}} & \tbc{\textbf{K5}} & \tbc{\textbf{K6}} \tabularnewline
    \midrule
    \tbc{QEC (kod powierzchniowy)} & \tbc{tak} & \tbc{cykle pomiaru syndromu} & \tbc{nie} & \tbc{dekoder czasu rzeczywistego} & \tbc{błąd operacji poniżej progu} & \tbc{bramka} \tabularnewline
    \tbc{ZNE} & \tbc{nie} & \tbc{przy powielaniu bramek} & \tbc{$m$ serii, wzrost wariancji} & \tbc{dopasowanie krzywej} & \tbc{szum skalowalny, stały, słaby} & \tbc{bramka} \tabularnewline
    \tbc{PEC} & \tbc{nie} & \tbc{wstawiane operacje} & \tbc{rosną wykładniczo} & \tbc{losowanie i~ważenie} & \tbc{znany, stały kanał każdej bramki} & \tbc{bramka} \tabularnewline
    \tbc{Mitygacja błędu odczytu} & \tbc{nie} & \tbc{nie} & \tbc{kalibracja, do $2^n$ obwodów} & \tbc{odwrócenie macierzy} & \tbc{klasyczny, stały błąd odczytu} & \tbc{bramka} \tabularnewline
    \tbc{DD} & \tbc{nie} & \tbc{impulsy w~przerwach} & \tbc{nie} & \tbc{brak} & \tbc{szum wolnozmienny} & \tbc{bramka} \tabularnewline
    \tbc{Redukcja (rozdz.~\ref{chapter:pdu},~\ref{chapter:dl_error_reduction})} & \tbc{nie} & \tbc{nie} & \tbc{kalibracja na zadaniach wzorcowych} & \tbc{model korekty} & \tbc{powtarzalny błąd zadania} & \tbc{problem} \tabularnewline
    \bottomrule
  \end{tabular}
\end{table}

% !TeX root = ./../main.tex

\section{Próbkowanie kwantowe jako model obliczeń}\label{sec:tb:qsamp}
\agentinfo{
Stan (PHD-TS-95, 2026-11-27): podrozdział skrócony i poprawiony po recenzji \agentref{eksperymenty/recenzje/rozdzial-2-sekcje-2.3-2.4.md}; usunięto podrozdział o operatorze zniekształcenia fazy (definicja, twierdzenie o odczycie), funkcję FCoSamp z postacią Fouriera i powtórzony wzór interferencji; lista cięć w \agentref{eksperymenty/recenzje/rozdzial-2-skracanie.md}.
QCoSamp i kodowanie danych w fazie pochodzą z pracy Wereszczyńskiego i in. \agentref{bib:wereszczynski_cosine_2020} (preprint) i nie są wkładem autora rozprawy.
Konwencja (za pracą źródłową i kodem z pracy magisterskiej \agentref{artykuły/magisterka_werner_final/chapters/03.tex}): argument x_D = -pi + suma po j od 0 do m-1 z pi d_j / 2^j, bit d_0 najbardziej znaczący, siatka od -pi do pi - pi/2^(m-1); wartość składnika jest prawdopodobieństwem wyniku 0 na kubicie odczytu.
Rozbieżności źródeł niewidoczne po skróceniu: warunek lambda^2 + gamma^2 <= 2 w pracy źródłowej (z tożsamości trygonometrycznej wynika <= 4) i wzór odtwarzania szeregu Fouriera dotyczyły usuniętego akapitu o funkcji FCoSamp; dziedzina [-1, 1] funkcji referencyjnej w artykule ICCS \agentref{artykuły/Quantum_Errors_ICCS/main.tex} należy do rozdziału 4.
Etykiety używane przez rozdział 4 (nie usuwać): subsec:tb:qsamp:sampling, subsec:tb:qsamp:qcosamp, def:tb:qsamp:sampling, def:tb:qsamp:component, def:tb:qsamp:argument, def:tb:qsamp:circuit, thm:tb:qsamp:readout, fig:tb:qsamp:component, eq:tb:qsamp:two-variables.
}

Metody redukcji błędów rozwijane w rozprawie korygują wynik obliczeń, który ma postać rozkładu prawdopodobieństwa, a nie pojedynczego ciągu bitów.
Podrozdział definiuje ten sposób prowadzenia obliczeń i metodę QCoSamp, na obwodach której opracowano metodę PDU (rozdział~\ref{chapter:pdu}); metoda ta i kodowanie danych w fazie pochodzą z pracy Wereszczyńskiego i~in.~\cite{wereszczynski_cosine_2020}, dostępnej jako preprint, i nie są wkładem autora rozprawy.

\subsection{Obliczenia oparte na próbkowaniu}\label{subsec:tb:qsamp:sampling}

Pojedyncze wykonanie obwodu kończy się pomiarem, który zwraca jeden stan bazowy (podrozdział~\ref{subsec:tb:measurement}); w \glslink{quantum-sampling}{próbkowaniu kwantowym} wynikiem obliczeń nie jest pojedynczy taki stan, lecz rozkład prawdopodobieństwa stanów bazowych, szacowany z wielu wykonań tego samego obwodu~\cite{wereszczynski_cosine_2020,werner_experiment_driven_2022}.

\begin{Definition}[Eksperyment próbkowania i znormalizowany histogram]\label{def:tb:qsamp:sampling}
Niech $U$ będzie operatorem ewolucji obwodu, a $\mathbb{E}$ bazą pomiarową złożoną ze stanów bazowych mierzonych kubitów.
\emph{Eksperymentem próbkowania} o liczbie powtórzeń $\nshots$ nazywa się $\nshots$-krotne przygotowanie stanu $U\ket{0\dots0}$ i jego pomiar w bazie $\mathbb{E}$, a \emph{znormalizowanym histogramem} eksperymentu funkcję $\tilde{\prob}(\ket{k})$ określoną na $\mathbb{E}$, której wartościami są częstości względne~\eqref{eq:tb01:estimator} wyników.
\end{Definition}

Oprócz obwodu i eksperymentu próbkowania obliczenie wymaga słownika, który stanom bazowym przypisuje znaczenie w dziedzinie problemu (wartość funkcji w zadanym punkcie, intensywność piksela, prawdopodobieństwo przynależności obiektu do klasy), oraz rekonstrukcji tych wielkości z histogramu~\cite{wereszczynski_cosine_2020}.

\subsection{Metoda QCoSamp}\label{subsec:tb:qsamp:qcosamp}

Metoda \gls{qcosamp}, czyli kwantowe próbkowanie szeregiem cosinusów, służy do budowy operatorów ewolucji, dla których prawdopodobieństwo wyniku pomiaru jednego kubitu jest ważoną sumą przesuniętych funkcji cosinus argumentu wprowadzonego do obwodu~\cite{wereszczynski_cosine_2020}.

\begin{Definition}[Składnik cosinusowy]\label{def:tb:qsamp:component}
\emph{Składnikiem cosinusowym} o częstotliwości (krotności argumentu) $n\in\{1,2,\dots\}$ i fazie $r\in\mathbb{R}$ nazywa się funkcję o wartościach w przedziale $[0,1]$
\begin{equation}\label{eq:tb:qsamp:component}
\coscomp_{n,r}(x)=\frac{1}{2}\bigl(1+\cos(nx+r)\bigr),\qquad x\in[-\pi,\pi).
\end{equation}
\end{Definition}

Wartość składnika powstaje przez \glslink{interference}{interferencję} na jednym kubicie: zgodnie z~\eqref{eq:tb01:interference} układ bramek $\mathrm{H}\,\mathrm{P}(\varphi)\,\mathrm{H}$ daje $\prob(\ket{0})=\coscomp_{n,r}(x)$ dla $\varphi=nx+r$, a argument $x$ wprowadza się do obwodu jako stan bazowy osobnego rejestru.
W pracy źródłowej składnikiem (ang.\ \emph{component}) nazywa się parę takich funkcji o wspólnej częstotliwości, a kubit, na którym zachodzi interferencja, kubitem pomocniczym (\emph{ancilla}); w rozprawie składnikiem jest pojedyncza funkcja, a kubit ten nazywa się kubitem odczytu.
\agenttask{
Pytanie do autora (recenzja 2.4-I1): czy zaakceptować odstępstwa od zapisu pracy źródłowej metody QCoSamp (składnik jako pojedyncza funkcja, nazwa kubit odczytu zamiast ancilla, oznaczenia własne składnika, obwodu i bloku fazowego) i czy potwierdzić u współautorów preprintu dwie usterki jego zapisu, których rozprawa po skróceniu już nie przytacza (warunek na współczynniki z prawą stroną 2 zamiast 4 oraz wzór odtwarzania bez odjęcia wyrazu stałego)?
}

\begin{Definition}[Kodowanie argumentu]\label{def:tb:qsamp:argument}
Niech $D=(d_0,\dots,d_{m-1})\in\{0,1\}^{m}$ będzie ciągiem bitów zapisanym w stanie bazowym $\ket{D}=\ket{d_0\dots d_{m-1}}$ \emph{rejestru argumentu} o $m$ kubitach.
Ciąg $D$ koduje argument
\begin{equation}\label{eq:tb:qsamp:argument}
x_D=-\pi+\sum_{j=0}^{m-1}\frac{\pi d_j}{2^{j}} ,
\end{equation}
a liczbę $2^{m}$ kodowanych wartości nazywa się \emph{rozdzielczością} argumentu.
\end{Definition}

Bit $d_0$ jest najbardziej znaczący, a kodowane argumenty tworzą siatkę równomierną o kroku $\Delta x=\pi/2^{m-1}$ od $-\pi$ do $\pi-\Delta x$.

\begin{Definition}[Obwód składnika]\label{def:tb:qsamp:circuit}
\emph{Obwodem składnika} $\coscomp_{n,r}$ nazywa się obwód o $m+1$ kubitach, złożony z rejestru argumentu i \emph{kubitu odczytu} $a$, jedynego kubitu podlegającego pomiarowi, inicjowany stanem $\ket{0}^{\otimes(m+1)}$ i opisany operatorem
\begin{equation}\label{eq:tb:qsamp:circuit}
\compcirc_{n,r}(D)=\mathrm{H}_a\,\compphase^{(a)}_{n,r}\,\mathrm{H}_a\,E(D),
\quad
\compphase^{(a)}_{n,r}=\Bigl[\,\prod_{j=0}^{m-1}\mathrm{CP}_{j\to a}\bigl(\tfrac{n\pi}{2^{j}}\bigr)\Bigr]\mathrm{P}_a(r-n\pi),
\quad
E(D)=\bigotimes_{j=0}^{m-1}\mathrm{X}^{d_j},
\end{equation}
gdzie $\mathrm{CP}_{j\to a}(\theta)$ oznacza bramkę przesunięcia fazy o kąt $\theta$ na kubicie $a$ kontrolowaną $j$-tym kubitem rejestru argumentu.
\end{Definition}

\begin{Theorem}[Odczyt składnika]\label{thm:tb:qsamp:readout}
Po ewolucji $\compcirc_{n,r}(D)\ket{0}^{\otimes(m+1)}$ prawdopodobieństwa wyników pomiaru kubitu odczytu wynoszą
\begin{equation}\label{eq:tb:qsamp:readout}
\prob(\ket{0}_a)=\coscomp_{n,r}(x_D),\qquad \prob(\ket{1}_a)=1-\coscomp_{n,r}(x_D).
\end{equation}
\end{Theorem}
\begin{proof}
Operator $E(D)$ przeprowadza rejestr argumentu w stan $\ket{D}$, a diagonalny blok $\compphase^{(a)}_{n,r}$ nie zmienia składowej $\ket{D}\ket{0}_a$ i zwiększa fazę składowej $\ket{D}\ket{1}_a$ o $r-n\pi+\sum_{j}n\pi d_j/2^{j}$, czyli o $nx_D+r$, więc teza wynika z~\eqref{eq:tb01:interference} dla $\varphi=nx_D+r$.
\end{proof}

\begin{figure}[htbp]
\centering
\begin{subfigure}[b]{0.58\textwidth}
\centering
% !TeX root = ./../../main.tex
% Obwód składnika cosinusowego QCoSamp dla rejestru argumentu o m = 4 kubitach.
% Rysunek własny (quantikz2); układ bramek zgodny z obwodem referencyjnym
% z pracy magisterskiej autora (artykuły/magisterka_werner_final/chapters/03.tex, listing lst:python_cos_phi).
% Obwód jest składany w pudełku i skalowany do szerokości otoczenia (\linewidth).
\ifdefined\tbqsampbox\else\newsavebox{\tbqsampbox}\fi
\begin{lrbox}{\tbqsampbox}
\begin{quantikz}[row sep=0.14cm, column sep=0.18cm, font=\small]
\lstick{$\ket{0}$} & \gate{\mathrm{X}^{d_0}}\gategroup[4,steps=1,style={dashed,rounded corners,inner sep=2pt},label style={yshift=1pt}]{$E(D)$} & & \gategroup[5,steps=5,style={dashed,rounded corners,inner sep=2pt},label style={yshift=1pt}]{$\compphase^{(a)}_{n,r}$} & \ctrl{4} & & & & & \\
\lstick{$\ket{0}$} & \gate{\mathrm{X}^{d_1}} & & & & \ctrl{3} & & & & \\
\lstick{$\ket{0}$} & \gate{\mathrm{X}^{d_2}} & & & & & \ctrl{2} & & & \\
\lstick{$\ket{0}$} & \gate{\mathrm{X}^{d_3}} & & & & & & \ctrl{1} & & \\
\lstick{$a\colon\ket{0}$} & & \gate{\mathrm{H}} & \gate{\mathrm{P}(r-n\pi)} & \gate{\mathrm{P}(n\pi)} & \gate{\mathrm{P}\bigl(\tfrac{n\pi}{2}\bigr)} & \gate{\mathrm{P}\bigl(\tfrac{n\pi}{4}\bigr)} & \gate{\mathrm{P}\bigl(\tfrac{n\pi}{8}\bigr)} & \gate{\mathrm{H}} & \meter{}
\end{quantikz}
\end{lrbox}
\resizebox{\linewidth}{!}{\usebox{\tbqsampbox}}

\caption{składnik $\coscomp_{n,r}$ dla $m=4$}
\label{fig:tb:qsamp:component}
\end{subfigure}\hfill
\begin{subfigure}[b]{0.39\textwidth}
\centering
% !TeX root = ./../../main.tex
% Schemat blokowy obwodu realizującego szereg dwuskładnikowy (1/2)(kappa_{1,0} + kappa_{2,0}).
% Rysunek własny (quantikz2); struktura zgodna z obwodem z pracy magisterskiej autora
% (artykuły/magisterka_werner_final/graf/Eksperymenty/wiele_komp/uklad.png).
% Obwód jest składany w pudełku i skalowany do szerokości otoczenia (\linewidth).
\ifdefined\tbqsampbox\else\newsavebox{\tbqsampbox}\fi
\begin{lrbox}{\tbqsampbox}
\begin{quantikz}[row sep=0.2cm, column sep=0.2cm, font=\small]
\lstick{$b\colon\ket{0}$} & \gate{\mathrm{H}} & \gate[2]{\compphase^{(b)}_{2,0}} & & \gate[3]{\bigl(\compphase^{(ab)}_{3,0}\bigr)^{\dagger}} & \ctrl{2} & & \\
\lstick{$\ket{D}$} & \qwbundle{m} & & \gate[2]{\compphase^{(a)}_{1,0}} & & & & \\
\lstick{$a\colon\ket{0}$} & \gate{\mathrm{H}} & & & & \targ{} & \gate{\mathrm{H}} & \meter{}
\end{quantikz}
\end{lrbox}
\resizebox{\linewidth}{!}{\usebox{\tbqsampbox}}

\caption{szereg $\tfrac{1}{2}(\coscomp_{1,0}+\coscomp_{2,0})$}
\label{fig:tb:qsamp:two-components}
\end{subfigure}
\caption{Obwody metody QCoSamp: (a)~obwód składnika cosinusowego (definicja~\ref{def:tb:qsamp:circuit}) z kątami bramek fazowych dla $n=1$ i $r=0$; (b)~schemat obwodu szeregu dwuskładnikowego, w którym blok $(\compphase^{(ab)}_{3,0})^{\dagger}$, dodatkowo kontrolowany kubitem $a$, usuwa iloczyn faz $\mathrm{e}^{\mathrm{i}3x_D}$ ze składowej $\ket{1}_a\ket{1}_b$.
Opracowanie własne na podstawie~\cite{wereszczynski_cosine_2020,werner_magisterka_2022}.}
\label{fig:tb:qsamp:circuits}
\end{figure}

Obwód składnika (rysunek~\ref{fig:tb:qsamp:component}) wymaga $m$ kontrolowanych bramek fazowych, a stablicowanie składnika na całej siatce $2^{m}$ eksperymentów próbkowania, po jednym dla każdego ciągu $D$.
W odróżnieniu od pojedynczej bramki $\mathrm{P}(nx+r)$, która daje tę samą wartość, angażuje on wiele kubitów i \glslink{phase-kickback}{przeniesienie fazy} z bramek dwukubitowych (podrozdział~\ref{subsec:tb:gates}), dlatego to on służy w rozprawie do badania błędów.
Dla $n=1$ i $r=0$ otrzymuje się funkcję $\coscomp_{1,0}(x)=\tfrac{1}{2}(1+\cos x)$ i obwód $\compcirc_{1,0}(D)$, które w metodzie PDU pełnią rolę funkcji referencyjnej i obwodu referencyjnego (podrozdział~\ref{sec:pdu:reference}).

\emph{Szeregiem wieloskładnikowym} nazywa się kombinację wypukłą składników cosinusowych, $\cosseries(x)=\sum_{c}w_c\,\coscomp_{n_c,r_c}(x)$, gdzie $w_c>0$ i $\sum_{c}w_c=1$.
W obwodzie realizującym szereg składnik wskazują stany bazowe $\ket{c}$ dodatkowych kubitów pomocniczych, wprowadzonych w superpozycję: jeżeli przed końcową bramką Hadamarda układ jest w stanie $\sum_{c}\sqrt{w_c}\,\ket{c}\otimes(\ket{0}+\mathrm{e}^{\mathrm{i}\varphi_c}\ket{1})_a/\sqrt{2}$, to ze względu na ortogonalność stanów $\ket{c}$ składniki nie interferują ze sobą i $\prob(\ket{0}_a)=\sum_{c}w_c(1+\cos\varphi_c)/2$, czyli $\cosseries(x_D)$ dla $\varphi_c=n_cx_D+r_c$; niepożądane iloczyny faz różnych składników usuwa się przez \glslink{uncomputation}{dekomputację}.
W obwodzie z rysunku~\ref{fig:tb:qsamp:two-components}, użytym w pracy magisterskiej autora~\cite{werner_magisterka_2022}, kubit odczytu $a$ otrzymuje fazę $x_D$, kubit pomocniczy $b$ fazę $2x_D$, a po dekomputacji stan kubitu $b$ wskazuje, którą z~dwóch faz niesie kubit odczytu, więc $\prob(\ket{0}_a)=\tfrac{1}{4}(1+\cos x_D)+\tfrac{1}{4}(1+\cos 2x_D)$.

Funkcję dwóch zmiennych otrzymuje się, dodając drugi rejestr argumentu, przechowujący ciąg $D'$, i kontrolując nim drugi blok fazowy, o częstotliwości $n'$ i zerowej fazie, na tym samym kubicie odczytu; fazy pochodzące od obu rejestrów sumują się i
\begin{equation}\label{eq:tb:qsamp:two-variables}
\prob(\ket{0}_a)=\frac{1}{2}\bigl(1+\cos(nx_D+n'y_{D'}+r)\bigr)=\coscomp_{n,\,r+n'y_{D'}}(x_D),
\end{equation}
gdzie $y_{D'}$ jest argumentem kodowanym przez $D'$ zgodnie z~\eqref{eq:tb:qsamp:argument}.
Konstrukcję tego typu autor zastosował do funkcji zapisanej w pracy~\cite{werner_magisterka_2022} jako $\tfrac{1}{2}(1+\cos(x+y))$ (podrozdział~\ref{subsec:pdu:results:two}); stałego przesunięcia fazy w tamtym obwodzie źródło nie podaje jednoznacznie.

Drugim sposobem wprowadzania danych, opisanym w pracy źródłowej jako kodowanie danych stałych, jest zapisanie wektora liczb w \glslink{local-phase}{fazach lokalnych} stanów bazowych operatorem diagonalnym~\cite{wereszczynski_cosine_2020}; opiera się na nim reprezentacja LPIQE, zdefiniowana wraz z odczytem tych faz w podrozdziale~\ref{subsec:tb:qimg:lpiqe}.
\agenttask{
Pytanie do autora (recenzja 2.4-I3): po skróceniu rozprawa nie używa nazwy operator zniekształcenia fazy (PDO), która pochodzi z koncepcji z 2022 r., a nie ze źródeł opublikowanych, i myliła się z nazwą PDU; czy nazwa ma wrócić (jedno zdanie w tym akapicie albo w definicji LPIQE)?
}

% !TeX root = ./../main.tex

\section{Kwantowe reprezentacje obrazów}\label{sec:tb:qimg}
\agentinfo{
Stan (PHD-TS-95, 2026-11-27): podrozdział poprawiony po recenzji \agentref{eksperymenty/recenzje/rozdzial-2-sekcje-2.5-2.6.md} i skrócony (liczba pomiarów NEQR, błąd próbkowania FRQI i jego porównanie z LPIQE, rysunek obwodu LPIQE, zdanie o miarach używanych w rozdziałach; reprezentacje FRQI i NEQR bez własnego nagłówka); dopisano współczynniki korelacji rangowej; lista zmian w \agentref{eksperymenty/recenzje/rozdzial-2-skracanie.md}.
Ustalenie autora (2026-11-26): podrozdział krótki, ograniczony do tego, czego potrzebują rozdziały 5 i 7; bez przeglądu rodziny reprezentacji i bez tabeli zbiorczej.
Notacja LPIQE uporządkowana względem artykułów KES i CoDIT (stan znormalizowany, unitarny operator rekonstrukcji, permutacja jako bramka X na kubicie pomocniczym) i zgodna z implementacją \agentref{eksperymenty/lpiqe-pdu-pix2pix/lpiqe_pdu/lpiqe.py}, w której domyślny jest pierwszy wariant wzoru odczytu.
Oszacowania błędu próbkowania (metoda delta) są wyprowadzeniem własnym; dla LPIQE zgodne z \agentref{eksperymenty/lpiqe-pdu-pix2pix/results/sampling_error.csv}.
Otwarte: koder i parametry porównania z kompresją stratną (teza T2 jest robocza) ustala rozdział 5 po decyzji autora.
Rozdział 5 powinien przy liczbach z artykułów zaznaczyć, że źródła nie precyzują sposobu liczenia odchylenia standardowego błędu (recenzja, ustalenie Q11), i wskazywać wariant wzoru odczytu przy każdym odwołaniu do eq:tb:qimg:lpiqe-pixel.
Rysunek obwodu LPIQE usunięto dla objętości (plik figures/chp_theoretical_background/tb05_lpiqe_circuit.tex został; ma usterki składu opisane w recenzji, ustalenie Q14); po poprawieniu pasuje do rozdziału 5, obok opisu potoku.
}

\Gls{qimr} to sposób zapisania obrazu w stanie rejestru kwantowego wraz z procedurą odtworzenia go z wyników pomiarów.
Nośnikiem intensywności piksela może być kąt obrotu dodatkowego kubitu, stan bazowy ciągu kubitów albo faza lokalna stanu bazowego~\cite{werner_data_loss_2023}; praca wykorzystuje po jednym przedstawicielu tych grup: FRQI, NEQR i LPIQE (inne reprezentacje: przegląd~\cite{cai_survey_2018}).

\subsection{Schemat kodowania i rekonstrukcji, reprezentacje FRQI i NEQR}\label{subsec:tb:qimg:scheme}

Obraz o wysokości $H$ i szerokości $W$ pikseli, $\img=[\hat{p}_{r,c}]_{H\times W}$, gdzie $\hat{p}_{r,c}\in[0,1]$ jest intensywnością piksela w wierszu $r$ i kolumnie $c$, wektoryzuje się wierszami: $p_j=\hat{p}_{r,c}$ dla $j=rW+c$.
Indeks piksela jest zapisany w rejestrze pozycji o $n=\operatorname{log}_2W+\operatorname{log}_2H$ kubitach, czyli o $M=2^n=WH$ stanach bazowych $\ket{j}$; zakłada się, że $W$ i $H$ są potęgami dwójki, jak w obrazach $8\times8$ i $16\times16$ z rozdziału~\ref{chapter:pdu_images}, a przy innych rozmiarach i obrazach barwnych wykorzystuje się część stanów (rozdział~\ref{chapter:learned_encoding}).

Przetwarzanie obrazu na komputerze kwantowym składa się z czterech kroków: (1)~kodowania, czyli przygotowania stanu rejestru zależnego od obrazu, (2)~opcjonalnego przetwarzania tego stanu, pomijanego w rozdziale~\ref{chapter:pdu_images}, (3)~$\nshots$-krotnego pomiaru w bazie obliczeniowej oraz (4)~klasycznej rekonstrukcji intensywności $\tilde{p}_j$ z oszacowanego rozkładu.
Obraz odtworzony różni się od oryginału wskutek niedokładności reprezentacji, błędu próbkowania i szumu urządzenia, a redukcja błędów (podrozdział~\ref{sec:tb:qec}) działa po kroku~3 albo~4, na oszacowanym rozkładzie albo na odtworzonych intensywnościach, i nie zmienia obwodu.
Według tego samego schematu działa kodowanie uczone z rozdziału~\ref{chapter:learned_encoding}, w którym sieć klasyczna wyznacza kąty bramek stałego, płytkiego obwodu, a intensywność piksela jest przeskalowanym prawdopodobieństwem stanu bazowego, wygładzonym uczonym splotem; jest to kodowanie przybliżone, stratne także bez szumu.

Reprezentacja FRQI (ang.\ \emph{Flexible Representation of Quantum Images})~\cite{frqi} zapisuje intensywność w kącie $\theta_j$ kubitu barwy splątanego z rejestrem pozycji, a NEQR (ang.\ \emph{Novel Enhanced Quantum Representation})~\cite{neqr} zapisuje $q$-bitową wartość szarości $f_j\in\{0,\dots,2^q-1\}$ w stanie bazowym $q$ kubitów:
\begin{equation}\label{eq:tb:qimg:frqi}
\ket{I_{\mathrm{FRQI}}}=\frac{1}{\sqrt{M}}\sum_{j=0}^{M-1}\bigl(\cos\theta_j\ket{0}+\sin\theta_j\ket{1}\bigr)\ket{j},
\qquad
\ket{I_{\mathrm{NEQR}}}=\frac{1}{\sqrt{M}}\sum_{j=0}^{M-1}\ket{f_j}\ket{j}.
\end{equation}
Reprezentacja FRQI zajmuje $n+1$ kubitów, a jej pomiar daje stany $\ket{0}\ket{j}$ i $\ket{1}\ket{j}$ z prawdopodobieństwami $\cos^2\theta_j/M$ i $\sin^2\theta_j/M$, więc przy liniowym odwzorowaniu $\theta_j=\frac{\pi}{2}p_j$ intensywność można odtworzyć z liczb zliczeń $n_{0j}$ i $n_{1j}$ jako $\tilde{p}_j=\frac{2}{\pi}\arccos\sqrt{n_{0j}/(n_{0j}+n_{1j})}$.
W reprezentacji NEQR pojedynczy pomiar zwraca losową pozycję $j$ wraz z dokładną wartością $f_j$, więc bez szumu jedynym błędem jest kwantyzacja, lecz potrzeba $q-1$ kubitów więcej niż w FRQI (14 zamiast 7 dla obrazu $8\times8$ o 256 poziomach szarości); w badaniach~\cite{werner_data_loss_2023} uruchomiono ją wyłącznie na symulatorze.

\subsection{Reprezentacja LPIQE}\label{subsec:tb:qimg:lpiqe}

Reprezentacja LPIQE (ang.\ \emph{Local Phase Image Quantum Encoding}) zapisuje intensywność w fazie lokalnej stanu bazowego rejestru pozycji (podrozdział~\ref{subsec:tb:qsamp:qcosamp}).
Nie jest ona wkładem autora rozprawy: kodowanie obrazu w fazach lokalnych pochodzi z pracy Wereszczyńskiego i~in.~\cite{wereszczynski_cosine_2020}, nazwa LPIQE pojawia się w pracy~\cite{werner_data_loss_2023}, a operator rekonstrukcji w pracy~\cite{werner_error_reduction_2025}; w obu tych pracach autor rozprawy badał metodę eksperymentalnie.

\begin{Definition}[LPIQE]\label{def:tb:qimg:lpiqe}
Operatorem kodowania LPIQE obrazu $\img$ jest operator diagonalny $\lpiqe(\img)=\operatorname{diag}(e^{\mathrm{i}p_0},\dots,e^{\mathrm{i}p_{M-1}})$ na rejestrze pozycji, w którym intensywność $p_j\in[0,1]$ jest kątem fazy w radianach, a reprezentacją LPIQE jest stan $\lpiqe(\img)\,\mathrm{H}^{\otimes n}\ket{0}^{\otimes n}=\frac{1}{\sqrt{M}}\sum_{j}e^{\mathrm{i}p_j}\ket{j}$.
Operatorem rekonstrukcji jest operator na $n+1$ kubitach
\begin{equation}\label{eq:tb:qimg:recon-op}
\reconst(\img)=(\identity\otimes\mathrm{X})(\identity\otimes\mathrm{H})\,\Theta(\img),
\qquad
\Theta(\img)=\lpiqe(\img)\otimes\ket{0}\bra{0}+\identity\otimes\ket{1}\bra{1},
\end{equation}
w którym ostatnim czynnikiem iloczynu tensorowego jest kubit pomocniczy $a$.
\end{Definition}

Sam stan LPIQE nie nadaje się do odczytu, bo w bazie obliczeniowej każdy stan $\ket{j}$ ma prawdopodobieństwo $1/M$ niezależnie od faz.
Operator rekonstrukcji, zgodnie z nazwą przyjętą w~\cite{werner_error_reduction_2025}, obejmuje dlatego także kodowanie: $\Theta(\img)$ zapisuje obraz od razu na $n+1$ kubitach, w gałęzi $a=0$, bramka Hadamarda nakłada ją na gałąź odniesienia, a bramka $\mathrm{X}$ zamienia etykiety wyników kubitu pomocniczego.
Względem zapisu w~\cite{werner_data_loss_2023,werner_error_reduction_2025} stan jest tu znormalizowany, operator $\Theta(\img)$ nie zawiera czynnika $1/\sqrt{2}$, z którym operator rekonstrukcji nie byłby unitarny, a permutację stanów zapisano jako bramkę $\mathrm{X}$.
Działając operatorem $\reconst(\img)$ na równomierną superpozycję $n+1$ kubitów, otrzymuje się $\prob(\ket{j0})=(1-\cos p_j)/(2M)$ i $\prob(\ket{j1})=(1+\cos p_j)/(2M)$, a stąd dwa warianty odtworzenia intensywności z częstości:
\begin{equation}\label{eq:tb:qimg:lpiqe-pixel}
\tilde{p}_j=\arccos\bigl(1-2M\,\tilde{\prob}(\ket{j0})\bigr)
\quad\text{lub}\quad
\tilde{p}_j=\arccos\bigl(2M\,\tilde{\prob}(\ket{j1})-1\bigr).
\end{equation}
Warianty te nie są równoważne: korzystają z częstości różnych stanów, więc przy skończonej liczbie pomiarów i w obecności szumu dają różne wyniki, a prace źródłowe nie podają, którego z nich użyto.
\agenttask{
Pytanie do autora (recenzja 2.5, ustalenie Q2): który wariant wzoru eq:tb:qimg:lpiqe-pixel liczył oryginalny kod eksperymentów z artykułów KES 2023 i CoDIT 2025 (częstość stanu z kubitem pomocniczym w stanie 0 czy 1)? Odtworzenie na symulatorze używa domyślnie wariantu pierwszego.
}
Obwód zajmuje $n+1$ kubitów, tyle samo co FRQI: 7 dla obrazu $8\times8$ i 9 dla $16\times16$; operator $\lpiqe(\img)$ zadawano w pracach źródłowych macierzą, a jego rozkład na bramki natywne pozostawiano transpilatorowi (podrozdział~\ref{subsec:tb:gates}).

Dla $p_j\in[0,1]$ bez zniekształceń $\prob(\ket{j0})\le(1-\cos1)/(2M)\approx0{,}23/M$, a pierwszy wariant pozostaje określony aż do częstości $1/M$, co daje zapas na zniekształcenie wyników przez szum~\cite{werner_error_reduction_2025}; wartości odtworzone leżą jednak w $[0,\pi]$, a nie w $[0,1]$, a argument spoza $[-1,1]$, możliwy w drugim wariancie, trzeba obciąć.
Czynnik $2M$ wzmacnia przy tym błąd estymacji częstości: z wariancji~\eqref{eq:tb01:estimator} i linearyzacji funkcji odczytu (wyprowadzenie własne) wynika dla pierwszego wariantu $\sigma(\tilde{p}_j)\approx\sqrt{M/\nshots}/\cos(p_j/2)$, czyli około 0,10 dla jasnych pikseli obrazu $8\times8$ przy $\nshots=8192$.
Oszacowanie to nie dotyczy drugiego wariantu, w którym w mianowniku stoi $\sin(p_j/2)$, więc błąd ciemnych pikseli jest wielokrotnie większy; wpływ szumu urządzenia, wykraczający poza błąd próbkowania, bada rozdział~\ref{chapter:pdu_images}.

\subsection{Miary jakości rekonstrukcji obrazu}\label{subsec:tb:qimg:metrics}

Niech $e_j=\tilde{p}_j-p_j$ będzie błędem $j$-tej z $N$ porównywanych intensywności z przedziału $[0,1]$ ($N=WH$ dla obrazu monochromatycznego i $3WH$ dla RGB), a $\bar{e}$ średnią tych błędów.
Błąd średniokwadratowy, średni błąd bezwzględny, odchylenie standardowe błędu i szczytowy stosunek sygnału do szumu definiuje się jako
\begin{equation}\label{eq:tb:qimg:errors}
\begin{gathered}
\mathrm{MSE}=\frac{1}{N}\sum_{j}e_j^2,
\qquad
\mathrm{MAE}=\frac{1}{N}\sum_{j}|e_j|,
\\
\sigma_e=\sqrt{\frac{1}{N}\sum_{j}(e_j-\bar{e})^2},
\qquad
\mathrm{PSNR}=10\lg\frac{1}{\mathrm{MSE}}\ \mathrm{dB},
\end{gathered}
\end{equation}
gdzie $\lg$ oznacza logarytm dziesiętny; ponieważ $\mathrm{MSE}=\bar{e}^{\,2}+\sigma_e^2$, MSE łączy obciążenie rekonstrukcji z rozrzutem błędu, dlatego podaje się obok niego $\sigma_e$.
MAE jest mniej wrażliwy na pojedyncze duże błędy, a PSNR wyraża MSE w skali logarytmicznej i jest miarodajny przy porównaniach na tych samych obrazach i przy tym samym rodzaju zniekształcenia~\cite{huynh-thu_scope_2008}.
Wartości odtworzone spoza $[0,1]$, możliwe w LPIQE, zmieniają te miary, więc rozdziały podają, czy przed ich obliczeniem stosuje się obcięcie.

Współczynnik korelacji Pearsona $r$ między $p_j$ i $\tilde{p}_j$ (nieokreślony, gdy jeden z obrazów jest stały) mierzy liniowy związek rekonstrukcji z oryginałem niezależnie od przesunięcia i skali intensywności, więc wysoka korelacja przy dużym MSE wskazuje na błąd systematyczny o charakterze afinicznym.
Współczynniki korelacji rangowej mierzą zależność monotoniczną i nie zmieniają się przy rosnących przekształceniach nieliniowych: $r_{\mathrm{S}}$ Spearmana jest współczynnikiem Pearsona rang obu ciągów, a $\tau$ Kendalla różnicą udziałów par obserwacji uporządkowanych zgodnie i niezgodnie~\cite{bib:statistics_wilc}.
P-wartości tych współczynników dotyczą hipotezy o braku korelacji i zakładają niezależne obserwacje; sąsiednie piksele są skorelowane, więc p-wartości liczone na pikselach są orientacyjne.

Miara SSIM~\cite{wang_image_2004} porównuje lokalne średnie, wariancje i kowariancję intensywności w odpowiadających sobie oknach $x$ i $y$ obu obrazów:
\begin{equation}\label{eq:tb:qimg:ssim}
\mathrm{SSIM}(x,y)=\frac{(2\mu_x\mu_y+C_1)(2\sigma_{xy}+C_2)}{(\mu_x^2+\mu_y^2+C_1)(\sigma_x^2+\sigma_y^2+C_2)},
\end{equation}
gdzie $C_1$ i $C_2$ są małymi stałymi stabilizującymi, a wynik dla obrazu jest średnią po oknach, równą 1 tylko dla obrazów identycznych; okno dobrane do obrazów o boku od 8 do 16 pikseli i wartości stałych podaje rozdział~\ref{chapter:learned_encoding}.

% !TeX root = ./../main.tex

\section{Uczenie maszynowe}\label{sec:tb:ml}
\agentinfo{
Stan (PHD-TS-95, 2026-11-27): podrozdział skrócony i poprawiony po recenzji \agentref{eksperymenty/recenzje/rozdzial-2-sekcje-2.5-2.6.md}; usunięto rysunek modelu pix2pix, miary klasyfikacji poza dokładnością, własności gry GAN, jądro kwantowe, wzór reguły przesunięcia parametru i przegląd sieci ekwiwariantnych; dopisano akapit o porównaniach sparowanych i wnioskowaniu statystycznym (etykieta par:tb:ml:stats) oraz gradient z modelu zastępczego; lista zmian w \agentref{eksperymenty/recenzje/rozdzial-2-skracanie.md}.
Zakres: wyłącznie pojęcia używane w rozdziałach 4 i 6--8 oraz w sekcjach 3.5 i 3.7; podrozdział 2.6.5 łączy dawne podrozdziały 2.1.6 i 2.6.5 (etykiety subsec:tb:lie, subsec:tb:ml:lie, eq:tb01:generator, eq:tb:ml:lie) i służy tylko akapitom par:lr:noise:generators i par:lr:cls:lie oraz zdaniu o kwantowych sieciach ekwiwariantnych w sekcji 3.7.
Architektury, hiperparametry i liczby z eksperymentów opisują rozdziały 6--8; tutaj tylko definicje.
Miary jakości obrazu i współczynniki korelacji definiuje podrozdział o reprezentacjach obrazów.
Notacja dla rozdziałów 6--8: generator i dyskryminator to makra netG i netD, sieć wyznaczająca kąty obwodu to Theta z indeksem w, liczba pomiarów to makro nshots.
Etykiety używane przez inne pliki (nie usuwać): etykiety podsekcji (w tym subsec:tb:lie), eq:tb:ml:metrics, eq:tb:ml:adam, fig:tb:ml:hybrid, eq:tb:ml:lie, eq:tb01:generator.
Akapit o wnioskowaniu statystycznym nie ma cytowań dla testu t, poprawki Holma i przekształcenia Fishera (brak zweryfikowanych wpisów w bazie); rozdziały 4, 7 i 8 nie odwołują się jeszcze do etykiety par:tb:ml:stats.
}

Podrozdział wprowadza tylko pojęcia używane w dalszych rozdziałach; podstawy wykłada monografia~\cite{goodfellow2016deep}.

\subsection{Uczenie nadzorowane i klasyfikacja obrazów}\label{subsec:tb:ml:supervised}

W uczeniu nadzorowanym dany jest zbiór danych $\dataset=\{(\mathbf{x}_i,\mathbf{y}_i)\}_{i}$, w którym każda próbka $\mathbf{x}_i$ (np. obraz jako wektor intensywności pikseli) ma przypisaną wartość docelową $\mathbf{y}_i$.
W klasyfikacji jest nią etykieta jednej z $K$ klas, w translacji obraz--obraz z rozdziału~\ref{chapter:dl_error_reduction} obraz odniesienia, a w autoenkoderze z rozdziału~\ref{chapter:learned_encoding} sam obraz wejściowy.
Model $f_{\mathbf{w}}$ o parametrach (wagach) $\mathbf{w}$ wyznacza predykcję $\hat{\mathbf{y}}=f_{\mathbf{w}}(\mathbf{x})$, a uczenie polega na minimalizacji średniej straty na zbiorze uczącym:
\begin{equation}\label{eq:tb:ml:risk}
\symloss(\mathbf{w})=\frac{1}{|\trainset|}\sum_{(\mathbf{x},\mathbf{y})\in\trainset}\ell\bigl(f_{\mathbf{w}}(\mathbf{x}),\mathbf{y}\bigr).
\end{equation}
Dla wartości ciągłych stratą $\ell$ jest średni błąd bezwzględny (MAE), nazywany wtedy stratą L1 i oznaczany $\ell_{\mathrm{L1}}$, albo błąd średniokwadratowy~\eqref{eq:tb:qimg:errors}.
Dla klasyfikatora zwracającego rozkład prawdopodobieństwa klas $\hat{\mathbf{y}}$ stratą jest entropia krzyżowa $-\sum_{k}y_k\ln\hat{y}_k$, w której $\mathbf{y}$ jest etykietą zapisaną jako wektor jedynkowy, a dla jednego wyjścia $\hat{y}\in(0,1)$ i etykiety $y\in\{0,1\}$ binarna entropia krzyżowa (BCE) $-y\ln\hat{y}-(1-y)\ln(1-\hat{y})$.

Dane dzieli się na rozłączne zbiory: uczący $\trainset$, walidacyjny $\valset$, na którym dobiera się hiperparametry i wybiera punkt kontrolny modelu (zapisany stan wag), oraz testowy $\testset$, używany jednokrotnie do oceny końcowej, ponieważ wybór najlepszego z wielu wariantów na podstawie wyniku na danym zbiorze zawyża ten wynik.
\Gls{overfitting} zachodzi, gdy model dopasowuje się do szczegółów i szumu zbioru uczącego: strata na $\trainset$ maleje, a na $\valset$ rośnie; ryzyko to jest tym większe, im większa jest pojemność modelu względem liczby próbek~\cite{goodfellow2016deep}.
Od przeuczenia trzeba odróżnić \glsdisp{data-leakage}{przeciek danych} (ang.~\emph{data leakage}): wpływ informacji ze zbioru ocenianego na uczenie lub wybór modelu, który zawyża ocenę niezależnie od rzeczywistej zdolności uogólniania~\cite{kapoor_leakage_2023}.
W rozdziałach~\ref{chapter:learned_encoding} i~\ref{chapter:classification} znaczenie mają trzy jego formy: te same lub niemal te same próbki w zbiorze uczącym i ocenianym, parametry przetwarzania wstępnego wyznaczone przed podziałem danych oraz dobór hiperparametrów lub punktu kontrolnego na zbiorze, na którym podaje się wynik.

Wynik klasyfikatora podsumowuje macierz pomyłek $\mathbf{C}$, której element $C_{kl}$ jest liczbą próbek klasy $k$ przypisanych do klasy $l$, i wyznaczana z niej dokładność
\begin{equation}\label{eq:tb:ml:metrics}
\mathrm{ACC}=\frac{\sum_{k}C_{kk}}{\sum_{k,l}C_{kl}}.
\end{equation}
Dokładność odnosi się do klasyfikatora trywialnego: przy $K$ równolicznych klasach zgadywanie daje $1/K$.

\paragraph{Porównania sparowane i wnioskowanie statystyczne.}\label{par:tb:ml:stats}
Dwie metody ocenione na tych samych jednostkach (obrazach, próbkach, ziarnach generatora liczb losowych) porównuje się parami, przez różnice wyników.
Przedział ufności średniej różnicy wyznacza się z rozkładu t~Studenta o liczbie stopni swobody o jeden mniejszej od liczby par (sparowany test t sprawdza wtedy hipotezę o zerowej średniej różnicy i zakłada jej rozkład normalny) albo metodą bootstrap~\cite{efron_bootstrap_1979}, czyli z rozkładu średnich obliczonych na próbach losowanych ze zwracaniem spośród zaobserwowanych różnic.
Testy nieparametryczne nie zakładają rozkładu normalnego~\cite{bib:statistics_wilc}: dokładny test znaków korzysta tylko ze znaków różnic, które przy braku różnicy między metodami mają rozkład dwumianowy z prawdopodobieństwem $1/2$, test Wilcoxona także z rang ich modułów, test permutacyjny wyznacza rozkład statystyki przez losowe zamiany metod w parach, a test McNemara~\cite{mcnemar_note_1947} porównuje dwa klasyfikatory na podstawie liczby próbek, które poprawnie rozpoznał tylko jeden z nich.
W rodzinie $k$ testów stosuje się poprawkę Holma: p-wartości uporządkowane rosnąco porównuje się kolejno z progami $\alpha/k$, $\alpha/(k-1)$, \dots, do pierwszej hipotezy, której nie można odrzucić.
Przedział ufności współczynnika korelacji $r$ z $N$ obserwacji wyznacza się przekształceniem Fishera ($\operatorname{artanh}r$ ma w przybliżeniu rozkład normalny o odchyleniu standardowym $1/\sqrt{N-3}$), a współczynnik determinacji $R^2$ jest udziałem wariancji danych wyjaśnionym przez dopasowany model.

\subsection{Sieci neuronowe i sieci splotowe}\label{subsec:tb:ml:cnn}

Neuron oblicza $h_j=a\bigl(\sum_{i}w_{ji}x_i+b_j\bigr)$, gdzie $w_{ji}$ są wagami, $b_j$ przesunięciem, a $a$ nieliniową funkcją aktywacji; warstwy neuronów złożone kolejno tworzą perceptron wielowarstwowy (MLP).
W pracy występują funkcje aktywacji $\mathrm{ReLU}(u)=\max(0,u)$, $\mathrm{LeakyReLU}$, sigmoidalna $1/(1+e^{-u})$, $\tanh$ oraz softmax, zamieniająca wektor na rozkład prawdopodobieństwa.
Gradient straty~\eqref{eq:tb:ml:risk} oblicza się propagacją wsteczną~\cite{rumelhart_learning_1986}, szczególnym przypadkiem automatycznego różniczkowania~\cite{goodfellow2016deep}, i szacuje na losowych paczkach (wsadach) próbek; przejście przez cały zbiór uczący nazywa się epoką.
Optymalizator Adam~\cite{kingma2014adam} aktualizuje wagi na podstawie wygładzonych wykładniczo estymat pierwszego i drugiego momentu gradientu $\mathbf{g}_t$:
\begin{equation}\label{eq:tb:ml:adam}
\mathbf{m}_t=\beta_1\mathbf{m}_{t-1}+(1-\beta_1)\mathbf{g}_t,\quad
\mathbf{v}_t=\beta_2\mathbf{v}_{t-1}+(1-\beta_2)\mathbf{g}_t^{2},\quad
\mathbf{w}_t=\mathbf{w}_{t-1}-\eta\,\frac{\hat{\mathbf{m}}_t}{\sqrt{\hat{\mathbf{v}}_t}+\epsilon},
\end{equation}
gdzie działania wykonuje się po składowych, daszek oznacza korektę obciążenia ($\hat{\mathbf{m}}_t=\mathbf{m}_t/(1-\beta_1^{t})$, podobnie $\hat{\mathbf{v}}_t$), $\eta$ jest współczynnikiem uczenia, $\beta_1,\beta_2\in[0,1)$ współczynnikami zaniku, a $\epsilon$ małą stałą.

W warstwie splotowej sieci splotowej (CNN)~\cite{lecun1998gradient} każdy neuron oblicza ważoną sumę wartości z okna wokół danego położenia, przy czym te same wagi (jądro splotu) są używane we wszystkich położeniach~\cite{goodfellow2016deep}; splot z krokiem 2 zmniejsza rozdzielczość mapy cech dwukrotnie, a splot transponowany ją zwiększa.
Normalizacja wsadowa~\cite{ioffe_batch_2015} standaryzuje aktywacje średnią i wariancją z bieżącej paczki, a dropout~\cite{srivastava_dropout_2014} losowo usuwa podczas uczenia neurony wraz z ich połączeniami, co ogranicza przeuczenie.
Sieć U-Net~\cite{ronneberger_u-net_2015} składa się ze ścieżki kodującej, w której rozdzielczość map cech maleje, i ścieżki dekodującej, do której poziomów dołącza się mapy cech z odpowiednich poziomów kodera (połączenia skrótowe)~\cite{isola_image_2017}.

\subsection{Generatywne sieci przeciwstawne}\label{subsec:tb:ml:gan}

Generatywna sieć przeciwstawna (GAN)~\cite{goodfellow2014generative} składa się z generatora $\netG$, który przekształca wektor losowy $\mathbf{z}$ w próbkę, i dyskryminatora $\netD$, który zwraca prawdopodobieństwo, że próbka pochodzi ze zbioru danych, a nie z generatora; w warunkowej sieci GAN obie sieci otrzymują ponadto informację warunkującą $\mathbf{x}$~\cite{mirza2014conditional}.
Sieci uczy się jednocześnie w grze minimaksowej z funkcją wartości
\begin{equation}\label{eq:tb:ml:gan}
V_{\mathrm{c}}(\netD,\netG)=\mathbb{E}_{\mathbf{x},\mathbf{y}}\bigl[\ln\netD(\mathbf{x},\mathbf{y})\bigr]+\mathbb{E}_{\mathbf{x},\mathbf{z}}\bigl[\ln\bigl(1-\netD(\mathbf{x},\netG(\mathbf{x},\mathbf{z}))\bigr)\bigr],
\end{equation}
którą dyskryminator maksymalizuje, czyli minimalizuje stratę BCE z etykietą $1$ dla próbek prawdziwych $\mathbf{y}$ i $0$ dla wygenerowanych.
Model pix2pix~\cite{isola_image_2017} stosuje taką sieć do translacji obraz--obraz: warunkiem jest obraz wejściowy $\mathbf{x}$, generator ma zwrócić obraz docelowy $\mathbf{y}$, a funkcja celu łączy składnik przeciwstawny ze stratą L1 o wadze $\lambda_{\mathrm{L1}}$:
\begin{equation}\label{eq:tb:ml:pix2pix}
\netG^{*}=\arg\min_{\netG}\max_{\netD}\Bigl(V_{\mathrm{c}}(\netD,\netG)+\lambda_{\mathrm{L1}}\,\mathbb{E}_{\mathbf{x},\mathbf{y},\mathbf{z}}\bigl[\ell_{\mathrm{L1}}\bigl(\netG(\mathbf{x},\mathbf{z}),\mathbf{y}\bigr)\bigr]\Bigr).
\end{equation}
Generatorem jest sieć U-Net, a dyskryminator PatchGAN jest siecią w pełni splotową, która ocenia osobno każdy fragment (łatę) obrazu; losowość wprowadza jedynie dropout aktywny również podczas wnioskowania, więc wyniki dla tego samego wejścia mogą się nieznacznie różnić między wywołaniami~\cite{isola_image_2017}.
W eksperymentach z rozdziału~\ref{chapter:dl_error_reduction} i preprintu~\cite{werner_deep_learning_manuscript} strata generatora jest sumą straty BCE dla wyjścia dyskryminatora i błędu MAE względem obrazu odniesienia z wagami $1$ i $100$, co odpowiada $\lambda_{\mathrm{L1}}=100$; obie sieci uczy się optymalizatorem Adam~\eqref{eq:tb:ml:adam}.

\subsection{Kwantowe uczenie maszynowe}\label{subsec:tb:ml:qml}

Kwantowe uczenie maszynowe (QML) obejmuje modele, w których część obliczeń wykonuje obwód kwantowy~\cite{schuld2019qml,benedetti2019vqc}; podstawowym z nich jest wariacyjny obwód kwantowy (VQC), czyli sparametryzowany obwód $U(\boldsymbol{\theta})$ o ustalonej strukturze (ansatz), którego parametry, np. kąty bramek obrotu, dobiera klasyczny optymalizator przez minimalizację funkcji kosztu~\cite{cerezo_variational_2021}.
Dane klasyczne koduje się w stanie kwantowym, np. w kątach obrotów~\cite{benedetti2019vqc}, a wyjściem modelu są wartości oczekiwane obserwabli albo prawdopodobieństwa stanów bazowych, szacowane na urządzeniu z $\nshots$ powtórzeń obwodu.

\Gls{hybrid-model} łączy obwód kwantowy z klasycznym przetwarzaniem wstępnym, które zamienia dane na parametry obwodu, i klasycznym przetwarzaniem końcowym wyników pomiaru; wszystkie te składniki mogą mieć uczone parametry~\cite{benedetti2019vqc,mitarai_quantum_2018}.
W rozdziałach~\ref{chapter:learned_encoding} i~\ref{chapter:classification} obwód nie ma własnych parametrów: sieć klasyczna $\Theta_{\mathbf{w}}$ wyznacza dla każdej próbki kąty $\boldsymbol{\theta}=\Theta_{\mathbf{w}}(\mathbf{x})$ obwodu o stałej strukturze, a wyjście modelu powstaje z wyników pomiaru przez klasyczne przetwarzanie końcowe (rys.~\ref{fig:tb:ml:hybrid}).
Uczone są więc wagi $\mathbf{w}$ sieci klasycznej i, jeśli występują, parametry przetwarzania końcowego (splot końcowy w rozdziale~\ref{chapter:learned_encoding}), a obwód jest ustaloną warstwą, nieliniową względem $\boldsymbol{\theta}$.
Gdy celem jest odtworzenie wejścia, model jest autoenkoderem hybrydowym: koderem jest sieć $\Theta_{\mathbf{w}}$ wraz z obwodem, reprezentacją pośrednią stan kwantowy, a dekoderem pomiar i przetwarzanie końcowe.

\begin{figure}[htb]
\centering
% Schemat hybrydowego modelu kwantowo-klasycznego z obwodem o stałej strukturze.
% Rysunek własny (PHD-TS-17); wstawiany przez \input w chapters/theoretical_background_06_ML.tex.
\begin{tikzpicture}[
  font=\small,
  >={Stealth[length=2mm]},
  box/.style={draw, rounded corners=2pt, align=center, minimum height=1.0cm, inner sep=3pt},
  flow/.style={->, thick},
  lbl/.style={font=\footnotesize, align=center},
]
% wejście i sieć klasyczna
\node[box, fill=gray!15, minimum width=0.8cm] (in) at (0,0) {$\mathbf{x}$};
\node[box, fill=blue!15] (net) at (1.9,0) {sieć\\klasyczna $\Theta_{\mathbf{w}}$};
\draw[flow] (in) -- (net);

% obwód o stałej strukturze
\foreach \y in {0.45,0,-0.45} {
  \node[font=\scriptsize] at (4.0,\y) {$\ket{0}$};
  \draw (4.25,\y) -- (7.35,\y);
}
\draw[fill=white] (4.6,-0.65) rectangle (5.6,0.65);
\node at (5.1,0) {$U(\boldsymbol{\theta})$};
\draw[dashed, fill=red!8] (5.85,-0.65) rectangle (6.45,0.65);
\node at (6.15,0) {$\channel$};
\foreach \y in {0.45,0,-0.45} {
  \draw[fill=white] (6.75,{\y-0.15}) rectangle (7.15,{\y+0.15});
  \draw (6.82,{\y-0.08}) arc (160:20:0.14);
  \draw[->, >={Stealth[length=1mm]}] (6.95,{\y-0.1}) -- (7.08,{\y+0.11});
}
\draw[dashed, rounded corners=3pt, gray!70!black] (3.7,-0.95) rectangle (7.5,0.95);
\node[lbl, text=gray!30!black] at (5.6,-1.2) {obwód o stałej strukturze};

% kąty z sieci klasycznej
\draw[flow] (net.north) |- (5.1,1.4) -- node[right, pos=0.4] {$\boldsymbol{\theta}=\Theta_{\mathbf{w}}(\mathbf{x})$} (5.1,0.65);

% odczyt, przetwarzanie końcowe, strata
\node[box, fill=blue!15, font=\footnotesize] (post) at (11.1,0) {przetwarzanie\\końcowe};
\draw[flow] (7.5,0) -- node[above, font=\scriptsize] {$\tilde{\prob}(\ket{b})$, $\langle O\rangle_{\boldsymbol{\theta}}$} node[below, font=\scriptsize] {$\nshots$ pomiarów} (post.west);
\node[box, fill=orange!20] (loss) at (13.45,0) {$\ell(\hat{\mathbf{y}},\mathbf{y})$};
\draw[flow] (post) -- node[above, font=\scriptsize] {$\hat{\mathbf{y}}$} (loss);

% gradient
\draw[->, dashed, thick, gray!60!black] (loss.south) -- ++(0,-1.0) -| (net.south);
\node[lbl, text=gray!30!black] at (6.9,-2.0) {gradient $\partial\symloss/\partial\mathbf{w}$: różniczkowanie automatyczne (symulator)\\lub model zastępczy urządzenia (dostrajanie na urządzeniu)};
\end{tikzpicture}

\caption[Hybrydowy model kwantowo-klasyczny z obwodem o stałej strukturze]{Hybrydowy model kwantowo-klasyczny z obwodem o stałej strukturze.
Sieć klasyczna $\Theta_{\mathbf{w}}$ wyznacza kąty $\boldsymbol{\theta}$ obwodu $U(\boldsymbol{\theta})$; blok $\channel$ oznacza szum (model szumu w symulatorze albo szum urządzenia).
Opracowanie własne.}
\label{fig:tb:ml:hybrid}
\end{figure}

Pochodne wyjścia obwodu względem $\boldsymbol{\theta}$, potrzebne do gradientu straty, wyznacza się w symulatorze automatycznym różniczkowaniem, a na urządzeniu można je oszacować regułą przesunięcia parametru, z dwóch wykonań obwodu z przesuniętym kątem na każdą składową gradientu~\cite{mitarai_quantum_2018,schuld_evaluating_2019}.
W rozdziałach~\ref{chapter:learned_encoding} i~\ref{chapter:classification} reguły tej się nie stosuje: w dostrajaniu na urządzeniu, czyli dalszym uczeniu nauczonego wcześniej modelu, wynik obwodu pochodzi z pomiaru, a gradient z różniczkowalnego modelu zastępczego urządzenia w symulatorze~\cite{wright_deep_2022}.
Uczenie z modelem szumu polega na zastąpieniu w symulatorze ewolucji unitarnej kanałami kwantowymi (podrozdział~\ref{subsec:tb:noise:models}); tak nauczony model jest dopasowany do założonego modelu szumu, a nie do urządzenia, którego szum może mieć inną strukturę i zmieniać się między kalibracjami.
Ponadto uczenie prowadzi się zwykle na dokładnych prawdopodobieństwach, a na urządzeniu dostępne są tylko estymaty obarczone szumem próbkowania; oba źródła błędu są w rozdziałach~\ref{chapter:learned_encoding} i~\ref{chapter:classification} analizowane osobno.

\subsection{Grupy Liego w bramkach kwantowych i sieciach neuronowych}\label{subsec:tb:ml:lie}\label{subsec:tb:lie}

Bramki jednokubitowe po pominięciu fazy globalnej należą do grupy Liego $\mathrm{SU}(2)$ macierzy unitarnych o wyznaczniku 1~\cite{hall2015lie}; elementy jej algebry Liego $\liealg{su}(2)$, czyli przestrzeni stycznej do grupy w elemencie neutralnym, zapisuje się jako $-\mathrm{i}\gen$ z hermitowskim, bezśladowym generatorem $\gen$, a odwzorowanie (mapa) wykładnicze przeprowadza je w grupę, co daje postać bramki
\begin{equation}\label{eq:tb01:generator}
  U(\theta)=\exp\Bigl(-\mathrm{i}\tfrac{\theta}{2}\gen\Bigr),
\end{equation}
w której obrotom~\eqref{eq:tb01:rotation} odpowiada $\gen=\sigma_a$.
Postać ta pozwala modelować błąd bramki zaburzeniem generatora, $\gen\mapsto\gen+\genpert$: zaburzenie hermitowskie zachowuje unitarność i odpowiada błędowi koherentnemu (odchyleniu osi lub kąta obrotu), a niehermitowskie daje operator nieunitarny~\cite{raubitzek_quantum_noise} (podrozdział~\ref{par:lr:noise:generators}).

Symetrię danych można wbudować w architekturę sieci~\cite{bronstein2021geometric}: odwzorowanie $f$ jest \glsdisp{equivariance}{ekwiwariantne} względem grupy $G$ działającej na wejściach przez przekształcenia $T_g$, a na cechach przez $T'_g$, gdy $f(T_g\mathbf{x})=T'_g f(\mathbf{x})$ dla wszystkich $g\in G$ i $\mathbf{x}$.
W preprincie Raubitzka i in.~\cite{raubitzek_lie_group_manuscript}, omawianym w podrozdziale~\ref{par:lr:cls:lie}, grupa Liego nie jest symetrią danych, lecz strukturą narzuconą na przekształcenie cech wewnątrz sieci: cechy osadza się w wektorach $\mathbf{h}_i$, na przykład dzieląc je na bloki, a warstwa działania grupy przekształca każdy z nich macierzą z grupy $G$, np. $\mathrm{SO}(3)$, $\mathrm{SU}(2)$ lub $\mathrm{SU}(3)$:
\begin{equation}\label{eq:tb:ml:lie}
\mathbf{h}_i\mapsto g_i\mathbf{h}_i,\qquad g_i=\exp\Bigl(\sum_{a=1}^{\dim\liealg{g}}c_{i,a}B_a\Bigr)\in G,
\end{equation}
gdzie $B_a$ są macierzami bazy algebry $\liealg{g}$ (w preprincie macierze bazy grup $\mathrm{SU}$ są hermitowskie, z jawnym czynnikiem $-\mathrm{i}/2$), a uczonymi parametrami są współrzędne $c_{i,a}$.
Parametryzacja przez algebrę gwarantuje, że $g_i$ należy do grupy przy dowolnych współrzędnych, a warstwę, klasyczny odpowiednik bramki obrotu~\eqref{eq:tb01:generator}, uczy się zwykłą propagacją wsteczną~\cite{raubitzek_lie_group_manuscript}.


\section{Podsumowanie}\label{sec:theoretical_background:summary}
\agentinfo{
Krótkie zebranie wprowadzonych pojęć z odesłaniem do rozdziałów, które z nich korzystają.
Pisać po ukończeniu wszystkich części rozdziału.
}
